\documentclass[11pt]{article}
% Shared preamble for manuscripts-v2 papers. Load with % Shared preamble for manuscripts-v2 papers. Load with % Shared preamble for manuscripts-v2 papers. Load with \input{../common/preamble.tex}
% immediately after \documentclass[11pt]{article}.
\ifdefined\pdftrailerid\pdftrailerid{}\fi
\usepackage[T1]{fontenc}
\usepackage{lmodern}
\usepackage[letterpaper,margin=1.02in]{geometry}
\usepackage{amsmath,amssymb,amsthm,mathtools,booktabs,array,microtype,enumitem}
\usepackage{tikz}
\usepackage{xurl}
\usepackage[colorlinks=true,allcolors=bluegreen]{hyperref}
\definecolor{bluegreen}{RGB}{22,83,91}
\usepackage[capitalize,nameinlink]{cleveref}

\theoremstyle{plain}
\newtheorem{theorem}{Theorem}[section]
\newtheorem{lemma}[theorem]{Lemma}
\newtheorem{proposition}[theorem]{Proposition}
\newtheorem{corollary}[theorem]{Corollary}
\newtheorem{conjecture}[theorem]{Conjecture}
\newtheorem*{theorem*}{Theorem}
\newtheorem{maintheorem}{Theorem}
\renewcommand{\themaintheorem}{\Alph{maintheorem}}
\theoremstyle{definition}
\newtheorem{definition}[theorem]{Definition}
\newtheorem{example}[theorem]{Example}
\newtheorem{question}[theorem]{Question}
\theoremstyle{remark}
\newtheorem{remark}[theorem]{Remark}
\crefname{maintheorem}{Theorem}{Theorems}
\crefname{conjecture}{Conjecture}{Conjectures}
\crefname{question}{Question}{Questions}
\crefname{example}{Example}{Examples}

\newcommand{\R}{\mathbb R}
\newcommand{\Q}{\mathbb Q}
\newcommand{\Z}{\mathbb Z}
\newcommand{\N}{\mathbb N}
\newcommand{\T}{\mathbb T}
\newcommand{\C}{\mathbb C}
\DeclareMathOperator{\den}{den}
\newcommand{\e}{\mathrm e}

% Immutable Lean source links at research-checkpoint-18.
\newcommand{\leancommit}{e4809cf561573aef588e6ffc491a35cdb79fd9ec}
\newcommand{\leanbase}{https://github.com/eldernos/MathematicsResearch/blob/\leancommit}
% \lean{File}{line}{label}: link to NormalityChecker/File.lean#Lline at the pinned commit.
\newcommand{\lean}[3]{\href{\leanbase/NormalityChecker/#1.lean\#L#2}{#3}}
% \repo{path}{label}: link to any repository path at the pinned commit.
\newcommand{\repo}[2]{\href{https://github.com/eldernos/MathematicsResearch/tree/\leancommit/#1}{#2}}

% Evidence labels (see papers/WAVE2-CONTRACT.md).
\newcommand{\evidence}[1]{\textup{[\emph{#1}]}}

\setlength{\emergencystretch}{1.5em}
\raggedbottom

% immediately after \documentclass[11pt]{article}.
\ifdefined\pdftrailerid\pdftrailerid{}\fi
\usepackage[T1]{fontenc}
\usepackage{lmodern}
\usepackage[letterpaper,margin=1.02in]{geometry}
\usepackage{amsmath,amssymb,amsthm,mathtools,booktabs,array,microtype,enumitem}
\usepackage{tikz}
\usepackage{xurl}
\usepackage[colorlinks=true,allcolors=bluegreen]{hyperref}
\definecolor{bluegreen}{RGB}{22,83,91}
\usepackage[capitalize,nameinlink]{cleveref}

\theoremstyle{plain}
\newtheorem{theorem}{Theorem}[section]
\newtheorem{lemma}[theorem]{Lemma}
\newtheorem{proposition}[theorem]{Proposition}
\newtheorem{corollary}[theorem]{Corollary}
\newtheorem{conjecture}[theorem]{Conjecture}
\newtheorem*{theorem*}{Theorem}
\newtheorem{maintheorem}{Theorem}
\renewcommand{\themaintheorem}{\Alph{maintheorem}}
\theoremstyle{definition}
\newtheorem{definition}[theorem]{Definition}
\newtheorem{example}[theorem]{Example}
\newtheorem{question}[theorem]{Question}
\theoremstyle{remark}
\newtheorem{remark}[theorem]{Remark}
\crefname{maintheorem}{Theorem}{Theorems}
\crefname{conjecture}{Conjecture}{Conjectures}
\crefname{question}{Question}{Questions}
\crefname{example}{Example}{Examples}

\newcommand{\R}{\mathbb R}
\newcommand{\Q}{\mathbb Q}
\newcommand{\Z}{\mathbb Z}
\newcommand{\N}{\mathbb N}
\newcommand{\T}{\mathbb T}
\newcommand{\C}{\mathbb C}
\DeclareMathOperator{\den}{den}
\newcommand{\e}{\mathrm e}

% Immutable Lean source links at research-checkpoint-18.
\newcommand{\leancommit}{e4809cf561573aef588e6ffc491a35cdb79fd9ec}
\newcommand{\leanbase}{https://github.com/eldernos/MathematicsResearch/blob/\leancommit}
% \lean{File}{line}{label}: link to NormalityChecker/File.lean#Lline at the pinned commit.
\newcommand{\lean}[3]{\href{\leanbase/NormalityChecker/#1.lean\#L#2}{#3}}
% \repo{path}{label}: link to any repository path at the pinned commit.
\newcommand{\repo}[2]{\href{https://github.com/eldernos/MathematicsResearch/tree/\leancommit/#1}{#2}}

% Evidence labels (see papers/WAVE2-CONTRACT.md).
\newcommand{\evidence}[1]{\textup{[\emph{#1}]}}

\setlength{\emergencystretch}{1.5em}
\raggedbottom

% immediately after \documentclass[11pt]{article}.
\ifdefined\pdftrailerid\pdftrailerid{}\fi
\usepackage[T1]{fontenc}
\usepackage{lmodern}
\usepackage[letterpaper,margin=1.02in]{geometry}
\usepackage{amsmath,amssymb,amsthm,mathtools,booktabs,array,microtype,enumitem}
\usepackage{tikz}
\usepackage{xurl}
\usepackage[colorlinks=true,allcolors=bluegreen]{hyperref}
\definecolor{bluegreen}{RGB}{22,83,91}
\usepackage[capitalize,nameinlink]{cleveref}

\theoremstyle{plain}
\newtheorem{theorem}{Theorem}[section]
\newtheorem{lemma}[theorem]{Lemma}
\newtheorem{proposition}[theorem]{Proposition}
\newtheorem{corollary}[theorem]{Corollary}
\newtheorem{conjecture}[theorem]{Conjecture}
\newtheorem*{theorem*}{Theorem}
\newtheorem{maintheorem}{Theorem}
\renewcommand{\themaintheorem}{\Alph{maintheorem}}
\theoremstyle{definition}
\newtheorem{definition}[theorem]{Definition}
\newtheorem{example}[theorem]{Example}
\newtheorem{question}[theorem]{Question}
\theoremstyle{remark}
\newtheorem{remark}[theorem]{Remark}
\crefname{maintheorem}{Theorem}{Theorems}
\crefname{conjecture}{Conjecture}{Conjectures}
\crefname{question}{Question}{Questions}
\crefname{example}{Example}{Examples}

\newcommand{\R}{\mathbb R}
\newcommand{\Q}{\mathbb Q}
\newcommand{\Z}{\mathbb Z}
\newcommand{\N}{\mathbb N}
\newcommand{\T}{\mathbb T}
\newcommand{\C}{\mathbb C}
\DeclareMathOperator{\den}{den}
\newcommand{\e}{\mathrm e}

% Immutable Lean source links at research-checkpoint-18.
\newcommand{\leancommit}{e4809cf561573aef588e6ffc491a35cdb79fd9ec}
\newcommand{\leanbase}{https://github.com/eldernos/MathematicsResearch/blob/\leancommit}
% \lean{File}{line}{label}: link to NormalityChecker/File.lean#Lline at the pinned commit.
\newcommand{\lean}[3]{\href{\leanbase/NormalityChecker/#1.lean\#L#2}{#3}}
% \repo{path}{label}: link to any repository path at the pinned commit.
\newcommand{\repo}[2]{\href{https://github.com/eldernos/MathematicsResearch/tree/\leancommit/#1}{#2}}

% Evidence labels (see papers/WAVE2-CONTRACT.md).
\newcommand{\evidence}[1]{\textup{[\emph{#1}]}}

\setlength{\emergencystretch}{1.5em}
\raggedbottom

% Author information approved by the repository owner on 2026-11-25:
%   name: Christopher Anderson; designation: Independent Researcher.
% No academic affiliation, ORCID or personal contact details have been supplied.
% Do not add any that the owner has not approved.
\newcommand{\OwnerAuthor}{Christopher Anderson}
\newcommand{\OwnerAffiliation}{Independent Researcher}
\newcommand{\OwnerContact}{Corrections and comments: \url{https://github.com/eldernos/MathematicsResearch/issues}}
\newcommand{\ReleaseLabel}{manuscripts-v2.0.0}
\newcommand{\LicenseLine}{\copyright\ 2026 Christopher Anderson. Licensed under \href{https://creativecommons.org/licenses/by/4.0/}{CC BY 4.0}.}
% \makebyline{version}{date}: centred author block used under the title.
\newcommand{\makebyline}[2]{%
  {\large \OwnerAuthor\par}%
  \ifx\OwnerAffiliation\empty\else{\small \OwnerAffiliation\par}\fi
  \vspace{0.3em}
  {\small Version #1 \enspace$\cdot$\enspace #2 \enspace$\cdot$\enspace \ReleaseLabel\par}%
  {\footnotesize \OwnerContact\par}%
  {\footnotesize \LicenseLine\par}%
  \hypersetup{pdfauthor={\OwnerAuthor}}}

\hypersetup{pdftitle={All-Prime Polynomial Normality and Optimal Irrationality Exponents}}
\newcommand{\dist}[1]{\lVert #1\rVert}
\newcommand{\fr}[1]{\{#1\}}
\DeclareMathOperator{\ord}{ord}
\DeclareMathOperator{\lcm}{lcm}
\DeclareMathOperator{\rad}{rad}
\newcommand{\ttt}[1]{\texttt{#1}}
\newcommand{\leanP}[3]{\leanat{\phaseXIXcommit}{#1}{#2}{#3}}
% Lean module ProtectedPole, pinned at the Phase 21 commit below.
\newcommand{\phaseXXIcommit}{347ec3fe631a3e7ab2ae25d53508a3a4c9ea9fa8}
\newcommand{\leanR}[3]{\leanat{\phaseXXIcommit}{#1}{#2}{#3}}
\renewcommand{\ReleaseLabel}{candidate draft, not yet released}
\setlength{\emergencystretch}{3em}

\begin{document}
\begin{center}
{\LARGE\bfseries All-Prime Polynomial Normality\\[3pt] and Optimal Irrationality Exponents\par}
\vspace{10pt}
\makebyline{1.0 (preprint draft)}{27 November 2026}
\end{center}

\begin{abstract}
Fix an integer base $b\ge2$ and an odd prime $p\nmid b$. Call a series $\sum_kc_kb^{-n_k}$ with rational
coefficients \emph{all-prime} if every prime not dividing $b$ eventually divides the reduced denominator of
every shifted prefix $b^n\sum_{n_k\le n}c_kb^{-n_k}$. We construct explicit all-prime series all of whose
nonconstant rational polynomial images are normal to base $b$, and we control their irrationality exponents.
A series with insertions at $2^{k^3}$ and denominators $p^{k^2}\ell_k$ has exact prefix denominators
$p^{k^2}\ell_1\cdots\ell_k$, and every nonconstant rational polynomial in it is normal and Liouville with
discrepancy $\exp(-c(\log N)^{2/3})$. On the polynomial schedules $tk^d$, auxiliary primes inserted at sparse
indices give all-prime series in several orders whose polynomial images of total degree $r<d$ are normal with
discrepancy $O((\log\log N)^{-d/(d-1)})$, provided the multiplicative order of $b$ modulo the $r$-th power of
the auxiliary modulus grows more slowly than $K^{d-r}$; we compute that order exactly. On geometric schedules
$t\lambda^k$ we obtain irrationality exponent exactly $\lambda$ in degree $r<\lambda$, and on the Stoneham
schedule exponent $p$ in all degrees. Finally, a doubling schedule with twisted numerators gives an all-prime
series all of whose nonconstant rational polynomial images are normal and whose irrationality exponent is
exactly $2$; a positive series with occupied positions $n_j\asymp j\sqrt{\log j}$, so that $n_{j+1}/n_j\to1$, has
the same three properties. The exponent $2$ comes from a determinant estimate for successive prefixes. By
contrast, distinct-prefix truncation separation can never certify an exponent below $(3+\sqrt5)/2$, and it
certifies a finite exponent exactly under an explicit gap condition on the positions.
\end{abstract}

\section{Introduction}\label{sec:intro}

Throughout, $b\ge2$ is an integer, $\T=\R/\Z$, $\fr{y}$ is the fractional part, $\dist{y}$ the distance to the
nearest integer and $e(\theta)=\exp(2\pi i\theta)$. For real $x$ and $N\ge1$ the \emph{extreme discrepancy} of
the radix orbit is
\[
D_N(x)=D_N(x;b)=\sup_{0\le a<c\le 1}\left|\frac{\#\{0\le n<N:\fr{b^nx}\in[a,c)\}}{N}-(c-a)\right|,
\]
and $x$ is normal to base $b$ if and only if $D_N(x)\to0$ \cite[Chapter~1, Theorem~8.1]{KN}. We write $v_p$ for
the $p$-adic valuation, $\ord_m(b)$ for the multiplicative order of $b$ modulo $m$, $\rad(b)$ for the product of
the primes dividing $b$, and $\mu(x)$ for the irrationality exponent of a real number $x$: the supremum of the
$\mu$ for which $|x-h/l|<l^{-\mu}$ has infinitely many solutions. Throughout, $p\nmid b$ is an odd prime,
$t=\ord_p(b)$, $a_0=v_p(b^t-1)$, and $\ell_1<\ell_2<\cdots$ are the primes not dividing $pb$.

For a series $x=\sum_kc_kb^{-n_k}$ with rational $c_k$ and strictly increasing positive integers $n_k$, the
\emph{shifted prefix} after $n$ shifts is $R_n=b^n\sum_{n_k\le n}c_kb^{-n_k}$, and similarly for a polynomial in such
series, expanded over tuples of insertions. The series is \emph{all-prime} if every prime not dividing $b$
divides the reduced denominator of $R_n$ for all large $n$. The orbit $\fr{b^nx}$ is close to the orbit of the
rational $R_n$ between insertions, and normality is decided by how fast those rationals equidistribute. Primes
dividing the denominators help, but only one pole needs to be protected; the others must not destroy the
arithmetic of the orbits. The theorems below show that every off-radix prime can be kept in the denominators
while all polynomial images are normal, and that this is compatible with irrationality exponent $2$.

\begin{maintheorem}[A universal Liouville example]\label{thm:A}
The series $W=\sum_{k\ge1}p^{-k^2}\ell_k^{-1}b^{-2^{k^3}}$ is all-prime: for $2^{k^3}\le n<2^{(k+1)^3}$ the reduced
denominator of $R_n$ is exactly $p^{k^2}\ell_1\cdots\ell_k$. For every nonconstant $P\in\Q[x]$ the number $P(W)$ is
normal to base $b$ and a Liouville number, and $D_N(P(W))\ll_{P,c}\exp(-c(\log N)^{2/3})$ for every
$c<\min\bigl(3(\log2)^{1/3},\,(\deg P)(\log p)(\log 2)^{-2/3}\bigr)$.
\end{maintheorem}

This is \Cref{thm:allprime}; \Cref{thm:protected} and \Cref{cor:dense} give general protected poles and denser
schedules. On polynomial schedules the consecutive positions have ratio tending to one, and the $p$-adic mass of
the shifted prefixes, studied in the companion manuscript \cite{PaperVIIv2}, takes over the role of the
protected pole.

\begin{maintheorem}[Polynomial schedules with growing auxiliary support]\label{thm:B}
Let $1\le s_1<\cdots<s_m$ and $r\ge1$ be integers, and $d\ge r+1$. Let $k_1<k_2<\cdots$ be indices prime to $p$,
$\mu_{k_j}=\ell_j$, $\mu_k=1$ otherwise, $\Lambda_K=\prod_{k_j\le K}\ell_j$, and
$W_a=\sum_{k\ge1}p^{-s_av_p(k)}\mu_k^{-1}b^{-tk^d}$. If
\begin{equation}\label{eq:PH}
\limsup_{K\to\infty}\frac{\log\ord_{\Lambda_K^r}(b)}{\log K}<d-r,
\end{equation}
then every $W_a$ is all-prime, and every nonconstant $P\in\Q[X_1,\ldots,X_m]$ of total degree at most $r$ satisfies
$D_N(P(W_1,\ldots,W_m))\ll_P(\log\log N)^{-d/(d-1)}$, and $D_N\ll_PN^{-\eta_P}$ if the largest weighted degree
$\max\sum_as_a\nu_a$ over its monomials is attained by a single monomial of degree one.
\end{maintheorem}

This is \Cref{thm:hybrid}. The order in \eqref{eq:PH} is computed exactly in \Cref{lem:exactorder}, and
\eqref{eq:PH} holds, for example, when $\log k_j\ge(\frac{r-1}{d-r}+\epsilon)j\log j$ for $r\ge2$. For these series we
can only show $|W_a-A/Q|\ge Q^{-\psi(Q)}$ with $\psi(Q)\to\infty$ (\Cref{prop:ratioone}); whether their
irrationality exponents are finite is open. Geometric schedules give exact finite exponents.

\begin{maintheorem}[Prescribed finite exponents]\label{thm:C}
\begin{enumerate}[label=\textup{(\roman*)},nosep]
\item For every $r\ge1$ with $p>(r+1)!$ and every integer $\lambda\ge\max(3,r+1)$ there is an all-prime series
$W_G=\sum_ku_kp^{-k}\mu_k^{-1}b^{-t\lambda^k}$, with positive integers $u_k,\mu_k$, such that $P(W_G)$ is normal to
base $b$ with $D_N(P(W_G))\ll_PN^{-\eta_P}$ for every nonconstant $P\in\Q[x]$ of degree at most $r$, and
$\mu(W_G)=\lambda$.
\item There is an all-prime series $W_S$ on the Stoneham schedule $p^k$ such that $P(W_S)$ is normal to base $b$
with $D_N(P(W_S))\ll_PN^{-\eta_P}$ for every nonconstant $P\in\Q[x]$, and $\mu(W_S)=p$.
\end{enumerate}
\end{maintheorem}

These are \Cref{thm:geometric,thm:stoneham}. Both exponents come from truncation separation, which compares an
approximation $A/Q$ with the prefixes of the series. We show that this certificate has an intrinsic limit.

\begin{maintheorem}[What truncation separation can certify]\label{thm:D}
Let a series have prefixes $U_k$ with $\log_b\den U_k=n_k+o(n_k)$ and $-\log_b(\xi-U_k)=n_{k+1}+o(n_{k+1})$, where $n_k$ are
strictly increasing positive integers with $\Lambda=\limsup n_{k+1}/n_k<\infty$. Distinct-prefix truncation
separation certifies a finite irrationality exponent if and only if
$\liminf_k\max_{i<k}(n_{i+1}-n_i)/n_k>0$, and then the certified exponent is at least
$\max(\Lambda,1+\Lambda/(\Lambda-1))\ge(3+\sqrt5)/2$, with equality for geometric positions of ratio $(3+\sqrt5)/2$.
\end{maintheorem}

This is \Cref{prop:certificate}. In particular it cannot certify exponent $2$, and it certifies nothing for
positions with ratio tending to one. Using the actual residues of successive prefixes instead
(\Cref{lem:determinant}) we reach the optimal exponent in all degrees.

\begin{maintheorem}[Exponent two in all degrees]\label{thm:E}
Let $\mu_k=\ell_j$ at the square indices $k=j^2$ and $\mu_k=1$ otherwise, and let $u_k$ be the least positive integer
with $u_k\equiv\mu_kb^{2^k}\pmod{p^k}$ and $u_k\equiv1\pmod{\rad(b)\mu_k}$. Put $W=\sum_{k\ge1}u_kp^{-k}\mu_k^{-1}b^{-2^k}$.
\begin{enumerate}[label=\textup{(\roman*)},nosep]
\item $W$ is all-prime, $\mu(W)=2$, $W$ is transcendental, and $D_N(P(W))\ll_PN^{-\eta_P}$ for every nonconstant
$P\in\Q[x]$.
\item Adding coefficients $1$ at suitable clusters of positions after each $2^k$ gives a series $W'$ with the same
properties except the rate (every $P(W')$ is normal), whose consecutive position ratios have $\liminf1$ and
$\limsup2$ and whose positions up to $N$ number $N^{1-o(1)}$.
\item Adding coefficients $1$ at all multiples of $h_k$ in $(2^k,2^{k+1}]$, where $h_k$ is the largest power of two at most
$\sqrt k$, gives a positive series $V$ whose occupied positions satisfy $n_j\asymp j\sqrt{\log j}$ and
$n_{j+1}/n_j\to1$, which is all-prime, satisfies $\mu(V)=2$, and has $D_N(P(V))\ll_PN^{-\eta_P}$ for every
nonconstant $P\in\Q[x]$.
\end{enumerate}
\end{maintheorem}

These are \Cref{thm:doubling,thm:dense}. Three mechanisms replace those of \Cref{thm:C}(ii): the twisted
numerators $u_k$ replace divisibility of the positions by $t$; an anchored upper-set lemma (\Cref{lem:upperset})
replaces the periodic pole retention of Stoneham numbers; and the determinant estimate replaces truncation
separation, which could certify at best $\max(2,3)=3$ at ratio $2$. For $V$ the dense positive support is a
feature of one expansion, while the arithmetic control comes from an exact regrouping into a second, lacunary
one. The construction does not determine the exponents of the polynomial-schedule series of \Cref{thm:B}: their
coefficients admit no such geometric regrouping.

\paragraph{Relation to earlier work.} Arithmetic constructions of normal numbers by sparse series go back to
Stoneham \cite{Stoneham73} and Korobov \cite{Korobov72}; Bailey and Crandall \cite[Theorem~4.8]{BC02} and Vandehey
\cite[Corollary~1.2 and Theorem~7.4]{Vandehey19} give normality criteria and power-saving estimates for incomplete
Korobov sums on moduli with \emph{fixed} prime support. The orbit estimate used here (\Cref{lem:uniformorbit})
is an elementary stride-differencing bound on prime-power moduli whose point is uniformity in the base
$b^\tau$ when $\tau$ grows with the auxiliary support; we claim no novelty for the method. Rings of numbers normal to
one base and not to another \cite{KS93} give normal polynomial images of suitable generators, and normal Liouville
numbers are known explicitly \cite{KS93,BC02}; absolutely normal Liouville numbers exist \cite{Bugeaud02}, and a
computable one was constructed in \cite{BHS15}. Polynomial images of Stoneham numbers are normal by
\cite[Theorem~A]{PaperIv2} and, for a related family, by Mira Cult's public note \cite[Theorem~1.1]{MiraPoly};
without auxiliary primes $W_S$ is the Stoneham number $\alpha_{b,p}$, and what \Cref{thm:C}(ii) adds is the growing
auxiliary support and the exponent. Manai \cite{Manai25} constructs explicit numbers $x$ such that $P(x)$ is normal,
indeed absolutely normal, for every nonconstant integer-valued polynomial $P$ (we rely on the abstract of that
preprint). Our theorems are of a different kind: one fixed base, every rational polynomial, explicit series whose
prefix denominators acquire every prime not dividing $b$, explicit rates, and a prescribed Liouville or finite
irrationality exponent. As far as we can tell neither result contains the other. Almost every real number is
normal and has irrationality exponent $2$, by the theorems of Borel and Khintchine, so the existence of normal
numbers with exponent $2$ is classical; we make no priority claim for explicit normal numbers with a prescribed
irrationality exponent as such, and we did not search that literature for this draft. The exponents in
\Cref{thm:C} come from the classical truncation argument for lacunary series (\Cref{lem:truncation}); the
determinant estimate of \Cref{lem:determinant} is in the spirit of the folding lemma for the continued fractions of
$\sum u^{-2^k}$ \cite{Shallit79,vdPS92}, adapted to corrections that introduce new primes. The companion manuscript
\cite{PaperVIIv2} studies the same polynomial schedules at the boundary $d=r$, where it proves, using a theorem of
Dietmann, Elsholtz, Kalmynin, Konyagin and Maynard \cite[Theorem~1]{DEKKM}, that $C(b,p,s,2)^2$ is irrational for
$s<(390/449)\,t\log b$; for the cube a conditional criterion involves gaps between sums of three cubes, for which
Ghidelli's theorem \cite[Theorem~1.1]{Ghidelli20} is the best available input. Those numbers have no auxiliary
primes and are not studied here. A detailed comparison, with the sources read, is kept in the ledger accompanying
the source files (\Cref{sec:formal}).

\paragraph{The companion manuscript.} We use from \cite{PaperVIIv2} (version~2.0, a candidate draft like this one)
the small-ball estimate for the limiting $p$-adic mass on the schedule $tk^d$ and its weighted analogue for several
orders, which we quote in \Cref{sec:hybrid} with their evidence labels. The short-orbit estimate and its splitting
corollary also appear there; we restate them with proofs in \Cref{sec:orbits}, together with the material on
growing auxiliary moduli that the companion does not need. Statement numbers of \cite{PaperVIIv2} refer to its
current draft.

\paragraph{Evidence labels.} Each numbered statement carries one label: \emph{Kernel-verified} (the complete
statement is a Lean declaration checked by the kernel, linked at a pinned commit), \emph{Lean-supported paper
proof} (named finite steps are Lean declarations), \emph{Paper proof}, or \emph{External input}. A statement quoted
from the companion manuscript carries the label it has there. A Lean-supported label means only that the named
finite step is formal: every exponential-sum estimate, every normality conclusion, every rate and every
Diophantine statement here is a written proof. Finite computations quoted as illustrations are measurements, not
evidence for asymptotic statements.

\section{Discrepancy, grids and exponential sums}\label{sec:prelim}

For a finite sample $u=(u_0,\ldots,u_{N-1})$ in $\T$ let $\mu_u$ be its empirical measure and
$D(u)=\sup_{0\le a<c\le1}|\mu_u([a,c))-(c-a)|$. A \emph{half-open arc} is the image in $\T$ of $[\alpha,\alpha+l)$ with
$0\le l\le1$, of length $|I|=l$.

\begin{lemma}[Arcs and comparison]\label{lem:compare}\evidence{Paper proof}
For every half-open arc $I$, $|\mu_u(I)-|I||\le D(u)$. If $u,v$ are samples of length $N$, $\epsilon\ge0$ and
$E=\#\{n<N:\dist{u_n-v_n}>\epsilon\}$, then $|D(u)-D(v)|\le E/N+2\epsilon$.
\end{lemma}

\begin{proof}
A wrapping arc is the complement of an interval, and its excess is the negative of that interval's excess. For
the comparison let $I=[\alpha,\alpha+l)$; if $l>1-2\epsilon$ then $\mu_u(I)-l<2\epsilon$, and otherwise every $n$ with $u_n\in I$
and $\dist{u_n-v_n}\le\epsilon$ has $v_n$ in an arc of length $l+2\epsilon$, so $\mu_u(I)\le l+2\epsilon+D(v)+E/N$. Exchange $u$
and $v$.
\end{proof}

\begin{lemma}[Block certificate]\label{lem:blocks}\evidence{Paper proof}
Let $x\in\R$, $N\ge1$, and let $B_1,\ldots,B_m\subset[0,N)$ be disjoint sets of integers with $U=N-\sum_i|B_i|$.
Suppose $y_n\in\T$ is given for $n\in\bigcup B_i$ and $|\#\{n\in B_i:y_n\in I\}-|B_i||I||\le H_i$ for every $i$ and every
half-open arc $I$. If $\dist{b^nx-y_n}\le\epsilon$ for all covered $n$ except $E$ of them, then
$D_N(x)\le(U+E+\sum_iH_i)/N+2\epsilon$.
\end{lemma}

\begin{proof}
Put $y_n=\fr{b^nx}$ at uncovered $n$; then $D(y)\le(U+\sum_iH_i)/N$, and \Cref{lem:compare} applies.
\end{proof}

We use the Erd\H{o}s--Tur\'an inequality \cite[Chapter~2, Theorem~2.5]{KN} in the form
\[
D(u)\le C\Bigl(\frac1{H+1}+\sum_{h=1}^H\frac1h\Bigl|\frac1N\sum_{n<N}e(hu_n)\Bigr|\Bigr)
\]
with an absolute constant $C$, and the lifting of exponents: if $\ell$ is an odd prime and $\ell\mid y-1$, then
$v_\ell(y^m-1)=v_\ell(y-1)+v_\ell(m)$ for $m\ge1$. We also use $\prod_{\ell\le x}\ell<4^x$ over primes (the primes in
$(m+1,2m+1]$ divide $\binom{2m+1}m\le4^m$) and Chebyshev's bound $\ell_k\ll k\log k$.

\begin{lemma}[Grids]\label{lem:grid}\evidence{Lean-supported paper proof}
Let $q=p^eQ$ with $\gcd(Q,pb)=1$, let $\tau\ge1$ satisfy $b^\tau\equiv1\pmod{pQ}$, and put $c=v_p(b^\tau-1)$. If $e\ge c$,
then for every integer $a$ prime to $p$, every $L\ge1$, every $L$ consecutive points $\fr{ab^j/q}$ and every half-open
arc $I$,
\[
\bigl|\#\{j:\fr{ab^j/q}\in I\}-L|I|\bigr|\le Lp^{c-e}+\tau p^{e-c}.
\]
\end{lemma}

\begin{proof}
Put $y=b^\tau$. By lifting the exponent $y$ has order $p^{e-c}$ modulo $p^e$ and generates the residues $\equiv1\pmod{p^c}$,
and $y\equiv1\pmod Q$, so by the Chinese remainder theorem $\{ay^m\bmod q\}=\{r:r\equiv a\pmod{p^cQ}\}$, a progression of
$p^{e-c}$ points with difference $p^{c-e}$ in $\T$. Splitting $j$ modulo $\tau$ gives $\tau$ such progressions per stretch of
length $\tau p^{e-c}$; a progression of $M$ points with difference $1/M$ meets each arc in $M|I|+\vartheta$ points with
$|\vartheta|\le1$. The order computation and the count over one stretch are the Lean declarations
\leanP{SurvivingGrid}{43}{\ttt{stride\_order}} and \leanP{PrimePowerGrid}{337}{\ttt{crt\_period\_count}}; the passage to
arbitrary $L$ is written.
\end{proof}

\section{Short orbits on growing moduli}\label{sec:orbits}

For an integer $G\equiv1\pmod p$ with $G\ne1$ put $A(G)=v_p(G-1)$.

\begin{lemma}[Uniform short orbits]\label{lem:uniformorbit}\evidence{Paper proof}
Let $\theta\in(0,1]$, $n_\theta=\lceil2/\theta\rceil+1$, $\delta=2^{-2-n_\theta}$ and $M_*=p^{n_\theta+8}$. Let $M_0\ge M_*$, let
$G\equiv1\pmod p$ be an integer with $G\ne1$ and $A=A(G)\le\frac3{16}\log_pM_0-1$, and let $e,c$ be integers with
$A<e\le\theta^{-1}\log_pM_0$ and $p\nmid c$. Then for every integer $0\le M\le M_0$,
\[
\bigl|S(G,e,c;M)\bigr|\le9M_0^{1-\delta},\qquad S(G,e,c;M)=\sum_{0\le m<M}e\Bigl(\frac{cG^m}{p^e}\Bigr).
\]
\end{lemma}

\begin{proof}
(i) By lifting the exponent, $v_p(G^n-1)=A+v_p(n)$ for $n\ge1$. (ii) Hence $G$ has order $p^{e-A}$ modulo $p^e$ and
generates $1+p^A\Z/p^e\Z$; for $X=e-A\ge1$ a complete period sums to $e(c/p^e)\sum_{z\bmod p^X}e(cz/p^X)=0$, so $|S|<p^X$ for
every $M$. (iii) For $H,w\ge1$, Cauchy--Schwarz applied to $H\cdot S=\sum_n\sum_{i<H}a_{n+iw}$ (summands extended by zero) gives
\[
|S|^2\le\frac{M+(H-1)w}{H}\Bigl(M+2\sum_{h=1}^{H-1}|C_h|\Bigr),\qquad
C_h=\sum_{0\le m<\max(0,M-hw)}e\Bigl(\frac{c(G^{m+hw}-G^m)}{p^e}\Bigr),
\]
and by (i) $c(G^{hw}-1)=c_hp^{A+v_p(hw)}$ with $p\nmid c_h$, so $C_h=S(G,e_h,c_h;\max(0,M-hw))$ with the \emph{same base} and
$e_h=e-A-v_p(hw)$.

Put $L_0=\log_pM_0\ge11$, $W=\lfloor\frac34L_0\rfloor$, $\kappa=\lfloor\frac12L_0\rfloor$, $H=\lfloor M_0^{1/4}\rfloor$, and let
$\mathcal F$ be the set of triples $(G,e,c)$ with $A(G)\le W/4$, $e-A(G)\ge1$, $p\nmid c$. We show by induction on $n$:
\emph{if $(G,e,c)\in\mathcal F$ and $X=e-A\le W+n\kappa$, then $|S(G,e,c;M)|\le C_nM_0^{1-\delta_n}$ for $0\le M\le M_0$}, with
$C_0=1$, $C_{n+1}=(4+8C_n)^{1/2}<9$ and $\delta_n=2^{-2-n}$. If $X\le W$, (ii) gives $|S|<p^W\le M_0^{3/4}$. If
$W+n\kappa<X\le W+(n+1)\kappa$, let $k$ be $X-A-W$ clipped to $[0,\kappa]$; then $X_1=X-A-k$ satisfies
$\frac34W\le X_1\le W+n\kappa$. Take $w=p^k\le M_0^{1/2}$. For $1\le h<H$, $j=v_p(h)\le\frac14L_0$ and the differenced triple has
excess $X_1-j\in[\frac5{16}L_0-\frac34,\,W+n\kappa]$, with $\frac5{16}L_0-\frac34\ge1$, so the case $n$ applies to every $C_h$. Since
$M+(H-1)w\le2M_0$ and $H\ge\frac12M_0^{1/4}$,
$|S|^2\le4M_0^{3/4}(M_0+2HC_nM_0^{1-\delta_n})\le(4+8C_n)M_0^{2-\min(1/4,\delta_n)}$, which is the case $n+1$. Finally
$A\le W/4$ and $e-A<L_0/\theta\le W+n_\theta\kappa$, because $L_0\ge n_\theta+8$; apply the case $n=n_\theta$.
\end{proof}

Some bound on $A(G)$ of order $\log_pM_0$ is necessary. For $G=1+p^a$ with $p^a\ge10$, $e=2a$, $c=1$ and
$M=\lfloor p^a/10\rfloor$, the geometric sum gives exactly
$|S|/M=\sin(\pi M/p^a)/(M\sin(\pi/p^a))\ge\sin(\pi/10)/(\pi/10)>0.983$, while $A=a=\log_p(10M)+O(1/M)$.

\begin{corollary}[Splitting off a coprime modulus]\label{cor:split}\evidence{Paper proof}
Let $y=a/(p^eQ)$ with $\gcd(a,pQ)=1$ and $\gcd(Q,pb)=1$, put $T_Q=\ord_Q(b)$, $\tau=\lcm(t,T_Q)$ and $G=b^\tau$. Then
$A(G)=a_0+v_p(T_Q)$, and for $T\ge1$ and $M_0=\lceil T/\tau\rceil$,
$|\sum_{j<T}e(b^jy)|\le\tau\max_{p\nmid c,\,M\le M_0}|S(G,e,c;M)|$. If \Cref{lem:uniformorbit} applies to $(M_0,G,e)$ and
$\tau\le T$, the sum is at most $18T(\tau/T)^\delta$. If $Q$ divides $\Lambda^rM'$, where $\Lambda$ is a product of primes not
dividing $pb$ and $\gcd(M',pb)=1$, then $A(G)\le a_0+v_p(\ord_{M'}b)+\max_{\ell\mid\Lambda}\log_p\ell$, an empty maximum being $0$.
\end{corollary}

\begin{proof}
Since $p\nmid t$, $\tau/t$ has the $p$-adic valuation of $T_Q$, and lifting the exponent from $b^t$ gives $A(G)$. Write
$y\equiv a_1/p^e+a_2/Q$ with $p\nmid a_1$; for $j=\tau m+j_0$, $e(b^jy)=e(b^{j_0}a_2/Q)\,e(b^{j_0}a_1G^m/p^e)$. Summing over the
at most $M_0$ values of $m$ gives the first bound, and $9\tau M_0^{1-\delta}\le18T^{1-\delta}\tau^\delta$. For the last claim, the
kernel of $(\Z/\ell^k)^\times\to(\Z/\ell)^\times$ is an $\ell$-group, so $v_p(\ord_{\ell^k}b)=v_p(\ord_\ell b)\le v_p(\ell-1)$.
\end{proof}

The fixed multiplier is harmless: $v_p(\ord_{M'}b)\le\log_pM'$ because $\ord_{M'}b\le M'$. The weak inequality is
needed, since equality holds at $M'=1$, where both sides of the valuation bound vanish.

\begin{lemma}[Largest auxiliary prime]\label{lem:largestprime}\evidence{Paper proof}
For $x\ge2$ let $A_b(x)=\lcm\{\ord_\ell b:\ \ell\le x\text{ prime},\ \ell\nmid pb\}$. Then $\log x\ll_b\log A_b(x)/\log\log A_b(x)$
for $x\ge x_0(b)$. Consequently, if $\Lambda$ is the product of all primes $\ell\le\ell_{\max}$ with $\ell\nmid pb$, then in
\Cref{cor:split} $A(G)\le a_0+v_p(\ord_{M'}b)+O_b\bigl(\log(\ord_\Lambda b+3)/\log\log(\ord_\Lambda b+3)\bigr)$.
\end{lemma}

\begin{proof}
Put $M=A_b(x)$. A prime with $\ord_\ell b=m$ divides $b^m-1$, which has at most $m\log_2b$ prime factors, so at most
$x^{2/3}\log_2b$ primes $\ell\le x$ have order at most $x^{1/3}$. Every other prime $\ell\le x$ with $\ell\nmid pb$ has order
$m>x^{1/3}$ with $m\mid M$ and $\ell\equiv1\pmod m$, and each such $m$ admits fewer than $x^{2/3}$ such primes. At most
$1+\omega(b)$ primes divide $pb$, so $\pi(x)-1-\omega(b)-x^{2/3}\log_2b\le d(M)x^{2/3}$ and $d(M)\ge x^{1/3}/(2\log x)$ for large
$x$. The divisor bound $\log d(M)\le(\log2+o(1))\log M/\log\log M$ \cite[Theorem~317]{HW08} gives the claim, and
$\ord_\Lambda b=A_b(\ell_{\max})$; bounded orders force a bounded $\ell_{\max}$ and are absorbed in the constant.
\end{proof}

\begin{lemma}[Orders below $x$]\label{lem:Ax}\evidence{Paper proof}
$\log\lcm_{\ell\le x}(\ell-1)\ll x\log\log x/\log x$, the lcm over primes $\ell$.
\end{lemma}

\begin{proof}
Prime powers $q^a$ dividing the lcm with $q\le x/\log x$ contribute at most $\pi(\sqrt x)\log x+\vartheta(x/\log x)\ll x/\log x$.
A prime $q>x/\log x$ divides it only if $mq+1$ is prime for some $m<\log x$, and the Selberg upper-bound sieve
\cite{HR74} bounds the number of $q\le x/m$ with $q$ and $mq+1$ both prime by $\ll\frac{x}{m(\log x)^2}\frac m{\varphi(m)}$; these
$q$ contribute $\ll\sum_{m<\log x}\frac{x}{\varphi(m)\log x}\ll\frac{x\log\log x}{\log x}$.
\end{proof}

The auxiliary budgets below involve $\ord_{\Lambda^r}(b)$, which can be computed exactly.

\begin{lemma}[Exact auxiliary orders]\label{lem:exactorder}\evidence{Paper proof}
Let $\Lambda$ be a squarefree product of primes not dividing $pb$, $r\ge1$, $O_\Lambda=\ord_\Lambda(b)$, and for odd $\ell\mid\Lambda$
let $w_\ell=v_\ell(b^{\ord_\ell b}-1)\ge1$.
\begin{enumerate}[label=\textup{(\roman*)},nosep]
\item If $\Lambda$ is odd, $\ord_{\Lambda^r}(b)=\lcm(O_\Lambda,E)$ with $E=\prod_{\ell\mid\Lambda}\ell^{\max(0,r-w_\ell)}$.
\item If $\Lambda=2\Lambda_o$ (so $b$ is odd), $\ord_{\Lambda^r}(b)=\lcm(\ord_{\Lambda_o^r}(b),2^{e_2(b,r)})$, where
\[
e_2(b,r)=\begin{cases}\max(0,\,r-v_2(b-1)),& b\equiv1\pmod4,\\ \max(1,\,r-v_2(b+1)),& b\equiv3\pmod4,\ r\ge2,\\ 0,& b\equiv3\pmod4,\ r=1.\end{cases}
\]
\item Consequently $\ord_{\Lambda^r}(b)\le\Lambda^{r-1}O_\Lambda\le\Lambda^{r-1}A_b(\max_{\ell\mid\Lambda}\ell)$; and if no odd $\ell\mid\Lambda$ is a base-$b$
Wieferich prime ($w_\ell=1$), then $\Lambda_o^{r-1}2^{e_2(b,r)}$ divides $\ord_{\Lambda^r}(b)$, with $\Lambda_o=\Lambda$ for odd $\Lambda$ and
$e_2=0$ then.
\end{enumerate}
\end{lemma}

\begin{proof}
For odd $\ell$, lifting the exponent from $b^{\ord_\ell b}$ gives $\ord_{\ell^r}b=\ord_\ell(b)\,\ell^{\max(0,r-w_\ell)}$. By the Chinese
remainder theorem $\ord_{\Lambda^r}b=\lcm_\ell\ord_{\ell^r}b$. For a prime $q$, $v_q$ of this lcm is the maximum over $\ell$ of
$v_q(\ord_\ell b)+[\ell=q]\max(0,r-w_q)$; since $v_q(\ord_qb)=0$ (the order divides $q-1$), this equals
$\max(v_q(O_\Lambda),\max(0,r-w_q))=v_q(\lcm(O_\Lambda,E))$, which proves (i). For (ii), $\ord_2b=1$; if $b\equiv1\pmod4$ then
$v_2(b^n-1)=v_2(b-1)+v_2(n)$, and if $b\equiv3\pmod4$ then $v_2(b^n-1)=1$ for odd $n$ and $v_2(b^n-1)=v_2(b+1)+v_2(n)$ for even
$n$; this gives $\ord_{2^r}b=2^{e_2(b,r)}$. Part (iii) follows from $\ord_{\ell^r}b\mid\ell^{r-1}\ord_\ell b$ and from (i)--(ii).
\end{proof}

The $2$-adic factor can be much smaller than $2^{r-1}$ even when $4\nmid b-1$: for $b=7$ and $r=4$, $7^2\equiv1\pmod{16}$, so
$\ord_{16}(7)=2<4$. The exact order permits earlier specific introductions of auxiliary primes; it gives no
improvement in the leading order $j\log j$ of the admissible positions for $r\ge2$, since every odd factor of $E$ must
still divide the stride when the full auxiliary power occurs.

\section{Protected poles}\label{sec:protected}

For a series $x=\sum_kc_kb^{-n_k}$ put $U_k=\sum_{j\le k}c_jb^{-n_j}$. For $P_1\in\Z[x]$ we write $R_n$ for the shifted
prefix of $P_1(x)$ and $q_n$ for its reduced denominator. If $q_n=p^eQ$ with $p\nmid Q$, $\tau=\lcm(t,\ord_Qb)$ and
$c=v_p(b^\tau-1)=a_0+v_p(\ord_Qb)\le e$, then by \Cref{lem:grid} any $T$ consecutive points of the orbit of $R_n$ meet
every half-open arc in $T|I|+\vartheta$ points with
\begin{equation}\label{eq:gridcount}
|\vartheta|\le Tp^{c-e}+\tau p^{e-c},\qquad \tau p^{e-c}\le t\,q_n .
\end{equation}

\begin{lemma}[Pole dominance]\label{lem:poledom}\evidence{Lean-supported paper proof}
Let $U,V\in\Q$ with $v_p(U)=-e'$ and $v_p(V)=-e$, where $0\le e'<e$. Let $P_1=\sum_{j\le r}a_jx^j\in\Z[x]$ have degree $r$ and
leading coefficient $A$, let $1\le l\le r$, and put
\[
F_l(U,V)=\sum_{i=0}^{l}\frac{P_1^{(i)}(U)}{i!}V^i=\sum_{j\le r}a_j\sum_{i\le\min(j,l)}\binom jiU^{j-i}V^i .
\]
If $\min(e',e-e')>v_p(A\binom rl)$, then $v_p(F_l(U,V))=v_p(A\binom rl)-((r-l)e'+le)$.
\end{lemma}

\begin{proof}
The term $U^{j-i}V^i$ has depth $D(j,i)=(j-i)e'+ie$. The term $(j,i)=(r,l)$ has valuation exactly $v_p(A\binom rl)-D(r,l)$;
every other term has an integer coefficient and valuation at least $-D(j,i)$, while
$D(r,l)-D(j,i)=(r-j)e'+(l-i)(e-e')\ge\min(e',e-e')$. The resulting divisibility $p^{le+(r-l)e'-g}\mid\den F_l(U,V)$,
$g=v_p(A\binom rl)$, is the Lean declaration \leanR{ProtectedPole}{347}{\ttt{truncated\_poly\_pole\_dvd\_den}}; it is a lower
bound for the depth, which is all that \eqref{eq:gridcount} needs.
\end{proof}

\subsection{A universal Liouville example}

\begin{theorem}[Universal all-prime normality]\label{thm:allprime}\evidence{Lean-supported paper proof}
Let $n_k=2^{k^3}$ and $W=\sum_{k\ge1}p^{-k^2}\ell_k^{-1}b^{-n_k}$.
\begin{enumerate}[label=\textup{(\roman*)},nosep]
\item For $n_k\le n<n_{k+1}$ the reduced denominator of $b^nU_k$ is exactly $p^{k^2}\ell_1\cdots\ell_k$.
\item For every nonconstant $P\in\Q[x]$ of degree $r$ the number $P(W)$ is normal to base $b$ and a Liouville number, and
$D_N(P(W))\ll_{P,c}\exp(-c(\log N)^{2/3})$ for every $c<c_0=\min\bigl(3(\log2)^{1/3},\,r\log p/(\log2)^{2/3}\bigr)$.
\item $W$ is transcendental, and for every $r$ the radix shifts of $(W,\ldots,W^r)$ are uniformly distributed in $\T^r$.
\end{enumerate}
\end{theorem}

\begin{proof}
(i) The number $p^{k^2}\ell_1\cdots\ell_k\,b^nU_k=\sum_{j\le k}p^{k^2-j^2}(\prod_{i\le k,i\ne j}\ell_i)b^{n-n_j}$ is an integer. Modulo $p$ only
the term $j=k$ survives, and modulo $\ell_j$ only the term $j$ survives; both are units. The exact $p$-power and the
divisibility by the newest prime $\ell_k$ are the Lean declarations
\leanR{ProtectedPole}{245}{\ttt{shifted\_prefix\_protected\_prime\_pow\_dvd\_den}} and
\leanR{ProtectedPole}{143}{\ttt{shifted\_prefix\_aux\_prime\_dvd\_den}}; the persistence of the earlier $\ell_j$ is the written
step just given.

(ii) Multiplying by $b^{e_0}$ changes $D_N$ by at most $e_0/N$, so write $b^{e_0}P(W)=P_1(W)/M'$ with $P_1=\sum_ja_jx^j\in\Z[x]$ of
leading coefficient $A$ and $M'$ prime to $b$; this changes $q_n$ by a divisor of $M'$, the $p$-depth by $v_p(M')$ and the
loss below by at most $\max_{\ell\mid M'}v_p(\ell-1)$. Put $V_k=p^{-k^2}\ell_k^{-1}b^{-n_k}$. Then $P_1(W)$ expands over ordered
tuples $(k_1,\ldots,k_j)$ at positions $n_{k_1}+\cdots+n_{k_j}$. Fix a shell $n_k\le n<n_{k+1}$ with $rn_{k-1}<n_k$ and
$(r+1)n_k\le n_{k+1}$, and discard the \emph{windows} $\mathcal W_{k,i}=[in_k,in_k+rn_{k-1}]$, $1\le i\le r$. A tuple with $i\ge1$ copies
of $k$ (and none larger) has position in $\mathcal W_{k,i}$, and tuples on smaller indices have position at most $rn_{k-1}<n_k$.
Outside the windows, therefore, $b^{-n}R_n=F_l(U_{k-1},V_k)$ with $l=\min(r,\lfloor n/n_k\rfloor)\ge1$, by Taylor's formula,
and the retained part of the shell consists of at most $r$ runs, on each of which $R_{n+1}=bR_n$.

Since $v_p(U_{k-1})=-(k-1)^2$, $v_p(V_k)=-k^2$ and $\min((k-1)^2,2k-1)>v_p(A)+\log_pr$ for $k\ge k_1(P)$, \Cref{lem:poledom}
gives at every retained shift
\[
v_p(q_n)=(r-l)(k-1)^2+lk^2-v_p\Bigl(A\binom rl\Bigr)\ge r(k-1)^2+2k-1-v_p(A)-\log_pr .
\]
Every present tuple has denominator dividing $p^{rk^2}\prod_{j\le k}\ell_j^r$, so $\log q_n\le rk^2\log p+O_r(k\log k)$ by
Chebyshev's bound and $\prod_{\ell\le x}\ell<4^x$. Writing $q_n=p^eQ_n$, the loss
$c=a_0+v_p(\ord_{Q_n}b)\le a_0+\log_p\ell_k=O(\log k)$ by \Cref{cor:split}, so \eqref{eq:gridcount} applies with
$e-c\ge r(k-1)^2+2k-1-O_P(\log k)$.

Put a buffer of length $H_b=\lceil2\log_b(N+2)\rceil$ before each window. Outside windows and buffers every absent tuple has
position larger than $n+H_b$; at most $(\log_2m)^{j/3}$ ordered $j$-tuples sit at position $m$ and all coefficients are at
most one, so $|b^nP_1(W)-R_n|\ll_P(\log N)^{r/3}b^{-H_b}\ll N^{-1}$. Let $K=\lfloor(\log_2N)^{1/3}\rfloor$ be the shell of $N$, and
apply \Cref{lem:blocks} with the retained runs as blocks: windows and buffers cost $O(n_{K-1}+K\log N)$ shifts; the at
most $(r+1)(K+1)$ runs cost $t\,q_n=e^{O_P(K^2)}$ each; and the mesh term $Tp^{c-e}$ costs
$O(n_{K-1}/N)+p^{-r(K-2)^2+O_P(\log K)}$. Since $n_{K-1}/n_K=2^{-3K^2+3K-1}$ and $K\ge(\log_2N)^{1/3}-1$, after division by $N$
these are $\exp(-3(\log2)^{1/3}(\log N)^{2/3}+O((\log N)^{1/3}))$ and $\exp(-r\log p(\log2)^{-2/3}(\log N)^{2/3}+O((\log N)^{1/3}))$,
which proves the rate.

For the Liouville property put $D_k=b^{n_k}p^{k^2}\ell_1\cdots\ell_k$, so $\log D_k=n_k\log b+O(k^2)$. The rational $P(U_k)$ has
denominator dividing $MD_k^r$, where $M$ is a common denominator of the coefficients of $P$, and
$0<|P(W)-P(U_k)|\le2\max_{[0,1]}|P'|\,b^{-n_{k+1}}$, the difference being nonzero because $P(W)$ is normal, hence
irrational. Since $n_{k+1}/n_k\to\infty$, for every $h$ this is eventually below $(MD_k^r)^{-h}$.

(iii) A polynomial relation $P(W)=0$ contradicts normality, and every nonzero integer combination of $W,\ldots,W^r$ is a
nonconstant integer polynomial in $W$; apply Weyl's criterion \cite[Chapter~1, Theorem~6.2]{KN}.
\end{proof}

In $P(W)$ with $\deg P\ge2$ only $p$ and, when $l=r$, the newest prime $\ell_k$ are controlled; the older primes enter
through sums that may cancel, and no persistence is claimed for them in $P(W)$.

\subsection{General protected poles and denser schedules}

\begin{theorem}[Protected poles]\label{thm:protected}\evidence{Lean-supported paper proof}
Let $n_k$ be strictly increasing positive integers and $e_k\ge1$ integers, let $m_k\ge1$ be integers prime to $pb$, and
let $u_k$ be integers with $p\nmid u_k$ and $0<|u_k|\le p^{e_k}m_k$. Put $c_k=u_k/(p^{e_k}m_k)$, $W=\sum_kc_kb^{-n_k}$,
$\Lambda_k=\lcm(m_1,\ldots,m_k)$ and $\lambda_k=\max_{\ell\mid\Lambda_k}v_p(\ord_\ell b)$. Assume
\[
\text{(S) }\frac{n_{k+1}}{n_k}\to\infty;\quad \text{(D) }e_k-e_{k-1}\to\infty;\quad
\text{(G) }e_k-\lambda_k\to\infty;\quad \text{(P$_r$) }k\,p^{re_k}\,\pi_k=o(n_k),
\]
where $\pi_k=\Lambda_k^r$ in general and $\pi_k=\Lambda_k^{r-1}A_b(\max_{\ell\mid\Lambda_k}\ell)$ if $\Lambda_k$ is squarefree. Then $P(W)$ is
normal to base $b$ and a Liouville number for every $P\in\Q[x]$ with $1\le\deg P\le r$, and, with $K$ maximal such that
$n_K\le N$,
\[
D_N(P(W))\ll_P\frac{n_{K-1}}{n_K}+p^{\lambda_{K-1}-e_{K-1}}+p^{\lambda_K-e_K}+\frac{Kp^{re_K}\pi_K}{n_K}+\frac{K\log(2N)}N .
\]
\end{theorem}

\begin{proof}
Repeat the proof of \Cref{thm:allprime} with $e_k$ in place of $k^2$. By (D) the deepest term of $U_{k-1}$ is eventually
unique and the gap $\min(e_{k-1},e_k-e_{k-1})$ of \Cref{lem:poledom} tends to infinity, so $v_p(q_n)\ge(r-1)e_{k-1}+e_k-O_P(1)$;
the loss is at most $a_0+\lambda_k+O_P(1)$, so by (G) the mesh term of \eqref{eq:gridcount} is $O(p^{\lambda_k-e_k})$. Under (S)
there are $O(rK)$ runs with $K=O(\log N)$. The term $\tau p^{e-c}$ of \eqref{eq:gridcount} is at most
$t\,p^{re_k}\ord_Qb\cdot O_P(1)$ with $Q\mid\Lambda_k^rM'$, and $\ord_Qb\ll\pi_k$, by \Cref{lem:exactorder}(iii) in the squarefree
case; (P$_r$) makes the residual cost $o(n_k)$. The tail is $O_P((\log N)^rb^{-H_b})$. The Liouville bound uses
$\log(p^{e_k}\Lambda_k)=o(n_k)$, which follows from (P$_1$): if $\pi_k=\Lambda_k$ this is immediate, and in the squarefree case
$\Lambda_k$ divides $b^{\ord_{\Lambda_k}b}-1$ with $\ord_{\Lambda_k}b\le\pi_k$, so $\log\Lambda_k<\pi_k\log b=o(n_k)$. (An empty maximum,
when $\Lambda_k=1$, is $0$, and $A_b$ of it is $1$.)
\end{proof}

\begin{corollary}[Denser all-prime schedules]\label{cor:dense}\evidence{Paper proof}
Take $m_k=\ell_k$, $u_k=1$ and $e_k=k\lceil\log\log(k+16)\rceil$ in \Cref{thm:protected}; then $W$ is all-prime.
\begin{enumerate}[label=\textup{(\alph*)},nosep]
\item If $n_k=\lceil\exp(k(\log\log(k+16))^2)\rceil$, then $W$ is normal and a Liouville number.
\item If $r\ge2$ and $n_k=\lceil\exp((r-1+\epsilon)k\log k)\rceil$ with $\epsilon>0$, then $P(W)$ is normal and Liouville for every
$P\in\Q[x]$ with $1\le\deg P\le r$.
\item If $\log n_k/(k\log k)\to\infty$ and $n_{k+1}/n_k\to\infty$, then $P(W)$ is normal and Liouville for every nonconstant
$P\in\Q[x]$.
\end{enumerate}
\end{corollary}

\begin{proof}
(D) and (G) hold, since $\lambda_k\le\log_p\ell_k=O(\log k)$. By the prime number theorem $\log\Lambda_k=\vartheta(\ell_k)\sim k\log k$, and
$A_b(x)$ divides $\lcm_{\ell\le x}(\ell-1)$, so \Cref{lem:Ax} gives $\log A_b(\ell_k)\ll k\log\log k$. In (a) the exponent of
(P$_1$) is $e_k\log p+\log A_b(\ell_k)=O(k\log\log k)=o(\log n_k)$; in (b) it is $(r-1)(1+o(1))k\log k+O(k\log\log k)<\log n_k$; in
(c) it is $o(\log n_k)$ for every $r$. (S) holds in all three cases, and the all-prime property is the argument of
\Cref{thm:allprime}(i).
\end{proof}

None of these schedules is shown to be optimal. In every complete-period construction of degree at least two the
$j$-th prime appears at a position $n$ with $\log n\gg j\log j$, because a run must contain a period carrying the factor
$\prod_{i\le j}\ell_i^{r-1}$.

\begin{remark}[Several generators]\label{rem:algebras}
If $\gamma_1,\ldots,\gamma_m>0$ are linearly independent over $\Q$, $f_k$ are reals with $f_k-f_{k-1}\to\infty$ and $f_k-f_{k-1}=o(f_k)$,
and $X_a=\sum_kc_{a,k}b^{-n_k}$ with $c_{a,k}$ as in \Cref{thm:protected} and depths $e_{a,k}=\lfloor\gamma_af_k\rfloor$, then under (S),
$\lambda_k=o(f_k)$ and $k(p^{\max_ae_{a,k}}\Lambda_k)^r=o(n_k)$ every nonconstant rational polynomial of total degree at most $r$ in
$X_1,\ldots,X_m$ is normal and Liouville: the maximal depth $\gamma\cdot\nu$ over the monomials is attained uniquely, and within it
the largest number of new occurrences wins by a multiple of $f_k-f_{k-1}$. We do not use this statement.
\end{remark}

\section{Polynomial schedules with growing auxiliary support}\label{sec:hybrid}

For integers $s,d\ge1$ put $C(b,p,s,d)=\sum_{i\ge1}p^{-sv_p(i)}b^{-ti^d}$. For $n\ge t$ let $K(n)=\lfloor(n/t)^{1/d}\rfloor$,
$J=J(n)=\lfloor\log_pK(n)\rfloor$ and $u=u(n)=n/(tp^{dJ})\in[1,p^d)$. Let $\Gamma$ be the set of positive rationals $a/p^h$ with
$p\nmid a$, with \emph{loss} $h$. For integers $1\le s_1<\cdots<s_m$ and a nonconstant $Y=\sum_\nu A_\nu X^\nu\in\Z[X_1,\ldots,X_m]$
of total degree $r$, let $S^*=\max\{\sum_as_a\nu_a:A_\nu\ne0,\nu\ne0\}$, let $T$ be the set of monomials attaining it, and let
\[
M_Y(u)=\sum_{\nu\in T}A_\nu\sum p^{\sum_{\rm slots}s_{a}h}\in\Z_p,
\]
where the inner sum runs over the ordered tuples of coordinates in $\Gamma$, one for each slot of $X^\nu$ (a slot carrying
the variable $a$ and a coordinate of loss $h$ has weight $p^{s_ah}$), whose $d$-th powers sum to at most $u$. This is the
\emph{weighted mass}; for $m=1$ and $Y=x^r$ it is the mass $M_r$ of \cite[\S7.5]{PaperVIIv2}. We quote two results of the
companion manuscript.

\begin{theorem}[Freezing {\cite[Lemma~7.10 and the proof of Theorem~7.20]{PaperVIIv2}}]\label{thm:freeze}\evidence{Paper proof}
Let $X_a=C(b,p,s_a,d)$ and $L\ge1$. If $J(n)\ge J_1(L)=\max(L,\lceil L/s_1\rceil+\lceil L/d\rceil)$, then the shifted prefix of
$Y(X_1,\ldots,X_m)$ is integral at every prime other than $p$, has $p$-adic valuation at least $-S^*J$, and
$p^{S^*J}b^{-n}R_n\equiv M_Y(u(n))\pmod{p^L}$.
\end{theorem}

\begin{theorem}[Weighted small balls {\cite[Theorems~7.18 and~7.24]{PaperVIIv2}}]\label{thm:wsmallball}\evidence{Paper proof}
Let $d\ge2$ and $r\le d-1$. For every $w\in\Z_p$ and $L\ge3$,
$|\{u\in[1,p^d):M_Y(u)\equiv w\pmod{p^L}\}|\ll_Y(\log L)^{-d/(d-1)}$. If $T$ consists of a single monomial of degree one,
$X_{a_*}$ say, then $M_Y(u)\equiv A_{a_*}\lfloor u^{1/d}\rfloor$ and the set is empty for $w=0$ and $L>v_p(A_{a_*})$.
\end{theorem}

For one order this is the small-ball theorem \cite[Theorem~7.18]{PaperVIIv2}, a Lean-supported paper proof there; for
several orders it is the small-ball estimate in the proof of \cite[Theorem~7.24]{PaperVIIv2}, a paper proof resting on
the weighted averaged isolation lemma \cite[Lemma~7.23]{PaperVIIv2}. The freezing threshold $J\ge L$ is needed when two
monomials differ in weighted degree by one.

\begin{theorem}[All-prime polynomial schedules]\label{thm:hybrid}\evidence{Paper proof}
Let $1\le s_1<\cdots<s_m$ and $r\ge1$ be integers, $d\ge r+1$, and let $k_1<k_2<\cdots$ be indices with $p\nmid k_j$. Put
$\mu_{k_j}=\ell_j$, $\mu_k=1$ otherwise, $\Lambda_K=\prod_{k_j\le K}\ell_j$, and
\[
W_a=\sum_{k\ge1}\frac{b^{-tk^d}}{p^{s_av_p(k)}\mu_k}\qquad(1\le a\le m).
\]
Assume \eqref{eq:PH}. Then:
\begin{enumerate}[label=\textup{(\roman*)},nosep]
\item for $K\ge p$ and $tK^d\le n<t(K+1)^d$ the reduced denominator of $b^n\sum_{k\le K}p^{-s_av_p(k)}\mu_k^{-1}b^{-tk^d}$ is
$p^{s_a\lfloor\log_pK\rfloor}\Lambda_K$; in particular each $W_a$ is all-prime;
\item for every nonconstant $P\in\Q[X_1,\ldots,X_m]$ of total degree at most $r$,
$D_N(P(W_1,\ldots,W_m))\ll_P(\log\log N)^{-d/(d-1)}$, and $D_N\ll_PN^{-\eta_P}$ if the largest weighted degree of $P$ is attained by
a single monomial of degree one.
\end{enumerate}
The hypothesis holds if $\log\Lambda_K=o(\log K)$, for example for $k_j=2^{2^{j-1}}$, and also if
$\log(tk_j^d)\ge(\frac{(r-1)d}{d-r}+\epsilon)j\log j$ eventually for $r\ge2$, or $\log(tk_j^d)\ge(\frac d{d-1}+\epsilon)\log A_b(\ell_j)$ for
$r=1$.
\end{theorem}

\begin{proof}
(i) Let $J=\lfloor\log_pK\rfloor\ge1$ and multiply by $p^{s_aJ}\Lambda_K$. Modulo $p$ only the indices $k=p^Jc$, $1\le c<p$, survive;
they are not auxiliary because $p\nmid k_j$, and contribute $\Lambda_Kb^{n-tk^d}\equiv\Lambda_Kb^n$ each, between $1$ and $p-1$ of them.
Modulo $\ell_j$ only the index $k_j$ survives, with an $\ell_j$-unit numerator.

(ii) Write $b^eP=Y/D$ with $Y$ integral and $D$ prime to $b$. At scale $J$ an auxiliary index has loss $J$, so in a monomial
of $T$ a slot holding it carries at least $p^{s_1J}\equiv0\pmod{p^L}$ after normalization once $J\ge L$; monomials outside $T$
vanish by the weight gap as in \Cref{thm:freeze}. (The hypothesis $p\nmid k_j$ is essential: an auxiliary prime at an index
of valuation $J$ would multiply loss-zero weights by $\ell^{-1}\not\equiv1$ and change the frozen residue.) Hence
\Cref{thm:freeze} holds for the $W_a$ with the same mass $M_Y$, and the shifted prefixes have denominators $p^{e'}Q'$ with
$Q'\mid D'\Lambda_K^r$, $D'$ the prime-to-$bp$ part of $D$, and $S^*J-L-O(1)\le e'\le S^*J+O(1)$ off the failed cells.

Take $L=\lfloor c_0\log N\rfloor$ with $c_0$ small as in \cite[Theorem~7.24]{PaperVIIv2}, and cut $[1,N]$ at the $O(N^{r/d})$
insertion positions, at the shell edges, at the $o(N^{r/d})$ cell edges of $M_Y\bmod p^L$ and at $\lceil\sqrt N\rceil$; every shift
$n\ge\sqrt N$ is frozen. Trim the final $B=\lceil2\log_bN\rceil$ positions of each run; off these buffers the real tail is $O(N^{-1})$,
because each coefficient at exponent $j$ is $O_P((j+1)^r)$ by tuple counting. By \eqref{eq:PH} choose $\epsilon_1>0$ with
$\log\ord_{\Lambda_K^r}b\le(d-r-5d\epsilon_1)\log K$ for large $K$, where $K=\lfloor(N/t)^{1/d}\rfloor$. Lifting the order from $\Lambda_K^r$ to
$D'\Lambda_K^r$ costs at most the fixed factor $D'\varphi(D')$, so the stride $\tau=\lcm(t,\ord_{Q'}b)$ is $O(N^{1-r/d-4\epsilon_1})$. Discard
the runs below $\lceil\sqrt N\rceil$, the runs with untrimmed length below $2N^{1-r/d-\epsilon_1}$, and the failed runs; by
\Cref{thm:wsmallball} and the cell count the failed runs cover $O(N(\log L)^{-d/(d-1)})$ shifts, and none if $T$ is a
single linear monomial. On a retained run of trimmed length $T'\ge N^{1-r/d-\epsilon_1}$, the Erd\H{o}s--Tur\'an inequality with
$\lceil N^\xi\rceil$ frequencies reduces to the orbit of $hR$ with reduced denominator $p^eQ$, $Q\mid D'\Lambda_K^r$ and
$c_1\log N\le e\le c_2\log N$ for a small fixed $\xi>0$, since $v_p(h)\le\xi\log_pN$. Here $M_0=\lceil T'/\tau\rceil\ge N^{2\epsilon_1}$. Since the
$\ell_j$ are introduced in increasing order, $\Lambda_K$ contains every prime $\ell\le\ell_{\max}$ with $\ell\nmid pb$, and
$\ord_{\Lambda_K}b\le\ord_{\Lambda_K^r}b\le N$; so \Cref{lem:largestprime} gives $A(b^\tau)=a_0+v_p(\ord_Qb)=O(\log N/\log\log N)$, which is
at most $\frac3{16}\log_pM_0-1$ and below $e$ for large $N$. With $\theta=\min(1,2\epsilon_1/(c_2\log p))$ we have $p^{\theta e}\le M_0$.
\Cref{cor:split} and \Cref{lem:uniformorbit} bound each harmonic by $18T'N^{-2\epsilon_1\delta}$. Summing over harmonics and
applying \Cref{lem:blocks} gives a power saving apart from the failed runs, which proves (ii), including $d=r+1$.

For the sufficient conditions: if $\log\Lambda_K=o(\log K)$ then $\ord_{\Lambda_K^r}b\le\Lambda_K^r=K^{o(1)}$. Otherwise
\Cref{lem:exactorder}(iii) gives $\log\ord_{\Lambda_K^r}b\le(r-1)\log\Lambda_K+\log A_b(\ell_{j(K)})$, where $\log\Lambda_{k_j}\sim j\log j$ by the prime
number theorem and $\log A_b(\ell_j)\ll j\log\log j$ by \Cref{lem:Ax}.
\end{proof}

For $m=1$ the statement is the all-prime polynomial schedule of degree at most $r$ in one order. The position
thresholds are sufficient and not proved necessary; the exact order of \Cref{lem:exactorder} may permit earlier
specific introductions.

\subsection{Irrationality exponents on polynomial schedules}

\begin{lemma}[Truncation separation]\label{lem:truncation}\evidence{Paper proof}
Let $\xi=\sum_kc_kb^{-n_k}$ with rationals $0<c_k\le C$ and strictly increasing integers $n_k$, and suppose $\den U_k$ divides
$D_kb^{n_k}$. If $8CQD_k\le b^{n_{k+1}-n_k}$, then either $A/Q=U_k$ or $|\xi-A/Q|\ge1/(2QD_kb^{n_k})$.
\end{lemma}

\begin{proof}
$0<\xi-U_k\le2Cb^{-n_{k+1}}\le1/(4QD_kb^{n_k})$, while $A/Q\ne U_k$ gives $|A/Q-U_k|\ge1/(QD_kb^{n_k})$.
\end{proof}

Put $Y_k=b^{n_{k+1}-n_k}/(8CD_k)$ and $\rho_k=n_k/(n_k-n_{k-1})$. For large $Q$ let $k$ be minimal with $Q\le Y_k$; then
$b^{n_k}<(8CD_{k-1}Q)^{\rho_k}$, so
\begin{equation}\label{eq:separation}
A/Q\ne U_k\ \Longrightarrow\ |\xi-A/Q|\ge\bigl(2D_k(8CD_{k-1})^{\rho_k}\bigr)^{-1}Q^{-1-\rho_k},
\end{equation}
while $A/Q=U_k$ forces $Q\ge\den U_k$ and $|\xi-U_k|\ge c_{k+1}b^{-n_{k+1}}$.

\begin{proposition}[Consecutive ratios tending to one]\label{prop:ratioone}\evidence{Paper proof}
For each series $W_a$ of \Cref{thm:hybrid} with $\log\Lambda_K=o(K)$ one has $|W_a-A/Q|\ge Q^{-\psi(Q)}$ for all large $Q$, with
$\psi(Q)=(1+o(1))\frac1d\bigl(\frac{\log_bQ}{td}\bigr)^{1/(d-1)}$.
\end{proposition}

\begin{proof}
Here $c_k\le1$, $D_k=p^{s_a\lfloor\log_pk\rfloor}\Lambda_k$, $\den U_k\ge b^{n_k}$, $\rho_k\sim k/d$ and $\log Y_{k-1}\sim tdk^{d-1}\log b$; substitute
into \eqref{eq:separation}. The case $A/Q=U_k$ gives exponent $n_{k+1}/n_k+o(1)\to1$.
\end{proof}

Since $\psi(Q)\to\infty$, this does not exclude Liouville approximations, and \Cref{prop:certificate} below shows that it is
the limit of truncation separation. \emph{Whether the irrationality exponent of $W_a$ is finite, and whether $W_a$ is a
Liouville number, is open.}

\begin{remark}[A quasi-polynomial schedule]\label{rem:quasipoly}
On the schedule $n_i=ti\lceil\exp((\log(i+1))^2)\rceil$, whose consecutive ratios also tend to one, inserting the $\ell_j$ at
indices prime to $p$ with $\log\Lambda_K=o((\log K)^2)$ into $\sum_ip^{-sv_p(i)}\mu_i^{-1}b^{-n_i}$ is expected to give an all-prime
series all of whose nonconstant rational polynomial images are normal, by combining a plateau estimate for the $p$-adic
mass on that schedule with \eqref{eq:gridcount}. We do not reproduce that argument and use it nowhere. For that series
\eqref{eq:separation} gives only $|W-A/Q|\ge Q^{-\psi(Q)}$ with $\psi(Q)=\exp((1+o(1))\sqrt{\log\log Q})$.
\end{remark}

\section{Truncation separation and finite exponents}\label{sec:finite}

\subsection{The limit of truncation separation}

Let $\xi$ be real and let $U_k<\xi$ be distinct increasing rational prefixes, $q_k=\den U_k$, $a_k=\log_bq_k$,
$t_k=-\log_b(\xi-U_k)$ and $g_k=t_k-a_k-\log_b2$. For $A/Q$ with $x=\log_bQ$, the gate $g_k\ge x$ is exactly
$\xi-U_k\le1/(2Qq_k)$, and then $A/Q\ne U_k$ gives $|\xi-A/Q|\ge1/(2Qq_k)$ by the triangle inequality. The \emph{distinct-prefix
certificate} at $Q$ uses the cheapest gated prefix and certifies the exponent
\[
\mu_{\rm dp}(x)=1+\frac{\inf_{g_k\ge x}a_k+\log_b2}{x};
\]
if $A/Q=U_k$ in reduced form, the actual exponent is $t_k/a_k$. The \emph{certified exponent}
$\mu_{\rm cert}=\max(\limsup_{x\to\infty}\mu_{\rm dp}(x),\limsup_kt_k/a_k)$ is an upper bound for $\mu(\xi)$, and the optimization over
$k$ is exact for this certificate, including arbitrary choices among several prefixes. No claim of optimality among
Diophantine methods is implied.

\begin{proposition}[Truncation certificates]\label{prop:certificate}\evidence{Paper proof}
Suppose $a_k=n_k+o(n_k)$ and $t_k=n_{k+1}+o(n_{k+1})$ for strictly increasing positive integers $n_k$ with
$\Lambda=\limsup_kn_{k+1}/n_k<\infty$.
\begin{enumerate}[label=\textup{(\roman*)},nosep]
\item $\mu_{\rm cert}<\infty$ if and only if
\begin{equation}\label{eq:gapcondition}
\liminf_{k\to\infty}\frac{\max_{i<k}(n_{i+1}-n_i)}{n_k}>0 .
\end{equation}
\item If $\mu_{\rm cert}<\infty$, then $\mu_{\rm cert}\ge\max\bigl(\Lambda,1+\Lambda/(\Lambda-1)\bigr)\ge(3+\sqrt5)/2$, and geometric positions with
limiting ratio $(3+\sqrt5)/2$ attain $\mu_{\rm cert}=(3+\sqrt5)/2$.
\end{enumerate}
In particular \eqref{eq:gapcondition} fails, and nothing finite is certified, when $n_{k+1}/n_k\to1$.
\end{proposition}

\begin{proof}
(i) The running maximum of $g_i$ over $i<k$ differs from $\max_{i<k}(n_{i+1}-n_i)$ by $o(n_k)$, uniformly after splitting off
finitely many early terms. If \eqref{eq:gapcondition} holds, the gates are unbounded, and at the first $k$ with $g_k\ge x$ one
has $x>\max_{i<k}g_i\gg n_k$, so $\mu_{\rm dp}(x)=O(1)$, while $t_k/a_k\to$ values at most $\Lambda$. Conversely, if
\eqref{eq:gapcondition} fails, choose $k\to\infty$ with $\max_{i<k}g_i=o(n_k)$, and $Q$ with $\log_bQ$ just above that maximum but
still $o(n_k)$ (for example also with $\log_bQ\ge\sqrt{n_k}$). Every gated prefix then has index at least $k$ and cost at least
$(1-o(1))n_k$, so $\mu_{\rm dp}\to\infty$ along these $Q$.

(ii) Here \eqref{eq:gapcondition} holds, so $\Lambda>1$. Take $x$ just above a record value $g_i$. Every gated prefix has index
above $i$ and cost at least $(1-o(1))n_{i+1}$, while $g_i\le(1-1/\Lambda+o(1))n_{i+1}$ because $n_i\ge(1/\Lambda-o(1))n_{i+1}$; so
$\mu_{\rm dp}\ge1+\Lambda/(\Lambda-1)-o(1)$ along the records. The equal-prefix exponents have $\limsup t_k/a_k=\Lambda$. The function
$\max(\Lambda,(2\Lambda-1)/(\Lambda-1))$ on $(1,\infty)$ is minimized where $\Lambda^2-3\Lambda+1=0$, at $\Lambda=(3+\sqrt5)/2$. For geometric
positions $g_i\sim(\Lambda-1)n_i$, and the bounds are attained. If $n_{k+1}/n_k\to1$, the preceding maximal gap is $o(n_k)$.
\end{proof}

Both outcomes of (i) occur with $\liminf n_{k+1}/n_k=1<\limsup n_{k+1}/n_k=2$: the positions $2^j,2^j+1$ satisfy
\eqref{eq:gapcondition}, while the concatenated integer intervals $[a_j,a_j^2]$, $a_1=2$, $a_{j+1}=2a_j^2$, violate it, since at
$a_j^2$ the preceding maximal gap is $O(a_j)$. The criterion concerns the positions, under the stated assumptions on the
denominators and tails.

\subsection{Geometric schedules and the Stoneham schedule}

\begin{theorem}[Geometric schedules]\label{thm:geometric}\evidence{Paper proof}
Let $r\ge1$, let $p\nmid b$ be a prime with $p>(r+1)!$, let $\lambda\ge r+1$ be an integer and $n_k=t\lambda^k$. Place $\ell_j$ at strictly
increasing indices $k_j\ge j$ with
\[
(\mathrm P_G)\qquad\limsup_{K\to\infty}\frac{\log\ord_{\Lambda_K^r}(b)}{K\log\lambda}<1,\qquad\Lambda_K=\prod_{k_j\le K}\ell_j .
\]
At $k=k_j$ let $\mu_k=\ell_j$ and let $u_k$ be the least positive integer with $u_k\equiv\ell_j\pmod p$ and $\gcd(u_k,b\ell_j)=1$; otherwise
$\mu_k=u_k=1$. Put $W_G=\sum_ku_kp^{-k}\mu_k^{-1}b^{-n_k}$. Then:
\begin{enumerate}[label=\textup{(\roman*)},nosep]
\item for $n_K\le n<n_{K+1}$ the reduced denominator of $b^nU_K$ is $p^K\Lambda_K$, and $W_G$ is all-prime;
\item for every nonconstant $P\in\Q[x]$ of degree at most $r$, $P(W_G)$ is normal to base $b$ and $D_N(P(W_G))\ll_PN^{-\eta_P}$;
\item $\lambda\le\mu(W_G)\le\max(\lambda,\frac{2\lambda-1}{\lambda-1})$; in particular $\mu(W_G)=\lambda$ when $\lambda\ge3$, and then $W_G$ is transcendental by
Roth's theorem \cite{Roth55}.
\end{enumerate}
The indices $k_j=j+\lceil2(r-1)j\log(j+1)/\log\lambda\rceil$ for $r\ge2$, and $k_j=j\lceil(\log\log(j+16))^2\rceil$ for $r=1$, satisfy
$(\mathrm P_G)$.
\end{theorem}

\begin{proof}
(i) The term at $K$ has $p$-adic valuation exactly $-K$, because $u_K\equiv\ell\not\equiv0$ or $u_K=1$, and earlier terms have larger
valuation; $\ell_j$ occurs only at $k_j$, with numerator prime to $\ell_j$.

(ii) \emph{Exact depth.} The hypotheses persist when $r$ is replaced by $\deg P$, so assume $\deg P=r$, and write
$b^{e_0}P(W_G)=Y(W_G)/D_0$ with $Y=\sum_aA_ax^a\in\Z[x]$, $A_r\ne0$, $\gcd(D_0,b)=1$. Write a degree-$a$ tuple as a sorted multiset
$\mathsf M=(i_1\ge\cdots\ge i_a)$ with position $\sum n_{i_c}$ and weight $w(\mathsf M)=\sum i_c$. If $\mathsf M<\mathsf M'$ lexicographically, first
differing at slot $c$, then the position of $\mathsf M$ is at most $\sum_{c'<c}n_{i'_{c'}}+(r-c+1)n_{i'_c-1}$, which is less than the
position of $\mathsf M'$ because $r<\lambda$. So the multisets present at $n$ form a lexicographic initial segment $\{\mathsf M\le g\}$, and
distinct multisets have distinct positions. If $\mathsf M<g$ first differs from $g$ at slot $c$, then
$w(\mathsf M)\le w(g_1,\ldots,g_{c-1},(g_c-1)^{r-c+1})$, with equality only for that candidate; so the top weight $w^*(n)$ is attained
by at most $r+1$ multisets, with between $1$ and $(r+1)!$ ordered tuples in all, and $r(K-1)\le w^*(n)\le rK$ on
$[n_K,n_{K+1})$. Every unit is $\equiv1\pmod p$ ($u/\mu\equiv1$, and $b^{-n_i}\equiv1$ since $t\mid n_i$), so the top tuples contribute
$p^{-w^*}b^n$ times a count prime to $p$, because $p>(r+1)!$, while other degree-$r$ tuples have valuation at least $-w^*+1$
and lower degrees have weights at most $aK\le w^*-((r-a)K-r)$. Hence $v_p(R_n)=-w^*(n)+v_p(A_r)-v_p(D_0)$ for all
$n\ge n_{K_0(P)}$, with no exceptional shifts.

\emph{Digits.} Let $K=K(N)\le\log_\lambda N$. There are $O((\log N)^r)$ insertion positions below $N$; the positions $n_k+(r-1)n_1$
are insertions of the top monomial with consecutive ratios at most $\lambda$, so the shifts before the first insertion after
$N^{1-\epsilon}$ number $O(N^{1-\epsilon})$; runs shorter than $N^{1-\epsilon}$ cover $O(N^{1-\epsilon}(\log N)^r)$ shifts and the buffers of length
$\lceil2\log_bN\rceil$ cover $O((\log N)^{r+1})$. On a remaining run starting at $n_0$, $K_0=K(n_0)\ge(1-\epsilon)\log_\lambda N-O(1)$, and
$q=p^vQ'$ with $v\in[r(K_0-1)-O(1),rK_0+O(1)]$ and $Q'\mid\Lambda_K^rD_0'$, where $D'_0$ is the prime-to-$p$ part of $D_0$. By
$(\mathrm P_G)$, $\ord_{\Lambda_K^r}b\le\lambda^{(1-4\epsilon)K}\le N^{1-4\epsilon}$ for small $\epsilon$, so $\tau\le N^{1-3\epsilon}$ and $M_0\ge N^\epsilon$;
$p^{\theta v}\le M_0$ with $\theta=\epsilon\log\lambda/(2r\log p)$; by \Cref{lem:largestprime} $A(G)=O(\log N/\log\log N)\le\frac3{16}\log_pM_0-1$;
and $e\ge v-\eta_0\log_pN>A(G)$ for small $\eta_0$. The Erd\H{o}s--Tur\'an inequality with $N^{\eta_0}$ frequencies, \Cref{cor:split}
and \Cref{lem:uniformorbit} give discrepancy $\ll N^{-\epsilon\delta}\log N+N^{-\eta_0}$ on each long run, and \Cref{lem:blocks}
completes the proof.

(iii) $\den U_k=D_kb^{n_k}$ exactly with $D_k=p^k\Lambda_k$, since the numerator is $\equiv u_k\Lambda_k/\mu_k$ modulo $b^{n_k-n_{k-1}}$ and
prime to $b$. The least admissible $u_{k_j}$ is at most $p\rad(b)\ell_j$, so $0<c_k\le C=\rad(b)$. Since
$\Lambda_k\mid b^{\ord_{\Lambda_k}b}-1$, $(\mathrm P_G)$ gives $\log\Lambda_k\le\lambda^{(1-\epsilon)k}\log b$, so $\log D_k=o(\lambda^k)$ and
$\log(1/c_{k+1})=o(\lambda^k)$. Lower bound: $|W_G-U_k|\le2Cb^{-n_{k+1}}=\den(U_k)^{-\lambda+o(1)}$. Upper bound: in
\eqref{eq:separation} $\rho_k=\lambda/(\lambda-1)$, so $A/Q\ne U_k$ gives exponent $\frac{2\lambda-1}{\lambda-1}+o(1)$, since $\log Q\gg n_k$ makes
the factors involving $D_k$ and $D_{k-1}$ equal to $Q^{o(1)}$; if $A/Q=U_k$, then $Q\ge D_kb^{n_k}$ and
$|W_G-U_k|\ge c_{k+1}b^{-\lambda n_k}\ge Q^{-\lambda-o(1)}$. Finally $\lambda\ge\frac{2\lambda-1}{\lambda-1}$ exactly when $\lambda\ge\frac{3+\sqrt5}2$.
\end{proof}

The hypothesis $p>(r+1)!$ is used only to prevent cancellation at ties, which are unavoidable with linear weights: for
$r=2$, $p=3$, $\lambda=3$ the top multisets $(7,7)$ and $(8,6)$ at $n=n_8+n_6$ carry $1+2=3$ orderings and the valuation
drops by one. Since $\lambda>r$ is needed for the lexicographic structure, all degrees with a finite exponent require a
different schedule.

\begin{remark}[Real ratios]\label{rem:realratio}
A real $\lambda>r$ works with $n_k=t\lceil\lambda^k\rceil$, once the sum starts at an index $k_0$ with $\lambda^{k_0}(\lambda-r)>2r$, so that the
lexicographic comparison stays strict; the rounded schedule has the same limiting ratios. Then $\mu(W_G)=\lambda$ for every
real $\lambda\ge\max(r+\epsilon,(3+\sqrt5)/2)$; in particular $\mu=(3+\sqrt5)/2$ is attained in degree $r\le2$.
\end{remark}

\begin{theorem}[The Stoneham schedule]\label{thm:stoneham}\evidence{Paper proof}
Place the $\ell_j$ at strictly increasing indices $k_j\ge j$ with $\log\Lambda_K=o(K)$, for example $k_j=\lceil j(\log(j+2))^2\rceil$. At
$k=k_j$ let $\mu_k=\ell_j$ and let $u_k$ be the least positive integer with $u_k\equiv\ell_j\pmod{p^j}$ and $\gcd(u_k,b\ell_j)=1$; otherwise
$\mu_k=u_k=1$. Put
\[
W_S=\sum_{k\ge1}\frac{u_k}{p^k\mu_k}\,b^{-p^k}.
\]
Then $W_S$ is all-prime, $P(W_S)$ is normal to base $b$ with $D_N(P(W_S))\ll_PN^{-\eta_P}$ for every nonconstant $P\in\Q[x]$, and
$\mu(W_S)=p$.
\end{theorem}

\begin{proof}
The all-prime property is as in \Cref{thm:geometric}(i), and the exponent is \Cref{thm:geometric}(iii) with $\lambda=p\ge3$,
$D_k=p^k\Lambda_k$ and $c_k\le p^{j-k_j}\rad(b)\le\rad(b)$; that part of the proof uses only the geometric ratio, so the absence of the
factor $t$ is harmless, and $\log D_k=o(p^k)$ follows from $\log\Lambda_K=o(K)$.

\emph{Pole retention.} Without auxiliary primes $W_S$ is the Stoneham number $\alpha_{b,p}$. For each $Y\in\Z[x]$ of degree $d$ there
is a precision $l$ such that $v_p(q_n)\ge dk-(l-1)-O_Y(1)$ on every late shell $p^k\le n<p^{k+1}$ of $\alpha_{b,p}$
\cite[Lemmas~2.2--2.4]{PaperIv2}; the residue of $p^{dk}R_d$ modulo $p^l$ that decides this involves only tuples whose reflected
losses $k-j_i$ total at most $l-1$, since the others carry $p^l$. In $W_S$ such a tuple uses only indices $j_i\ge k-l+1$, and
once $k\ge k_{l-1}+l$ every auxiliary index $k_J$ in that window has $J\ge l$, so its unit satisfies $u/\ell\equiv1\pmod{p^l}$, also when an
index is repeated. The residues modulo $p^l$ of $p^{dk}R_d$ for $W_S$ and $\alpha_{b,p}$ therefore coincide on late shells, and the
retention transfers. Lower degrees $a<d$ have $p^{ak}R_a$ integral at $p$, so in $p^{dk}R_Y$ they carry $p^{(d-a)k}$ and cannot
cancel the top degree once $k>l-1+v_p(A_d)$.

\emph{Digits.} Insertions number $O((\log N)^d)$, and at most $(\log_pm+1)^d$ tuples sit at position $m$, which a buffer of length
$\lceil2\log_bN\rceil$ absorbs. Run the argument of \Cref{thm:geometric} with $v\in[dk-O(1),dk+O(1)]$, $Q'\mid\Lambda_k^dD_0'$ and
$\log\ord_{\Lambda_k^d}b\le(d-1)\log\Lambda_k+\log A_b=o(k)$; here $\Lambda_k=N^{o(1)}$, so \Cref{cor:split} gives $A(G)=O(\log\log N)$ directly.
\end{proof}

The input from \cite{PaperIv2} consists of written proofs there (its Lemmas~2.2--2.4); only other statements of that
paper are Lean-supported or kernel-verified.

\section{Exponent two in all degrees}\label{sec:two}

\subsection{A determinant estimate}

For a reduced rational $u=a/q$, $q\ge1$, define
\[
B(u)=\inf_{l\ge1,\ h/l\ne u}l^2\Bigl|u-\frac hl\Bigr|=\min\Bigl\{\frac{|lL|}q:\ l,L\in\Z\setminus\{0\},\ L\equiv al\pmod q\Bigr\}>0 .
\]
Indeed $l^2|a/q-h/l|=l|al-qh|/q$, and $L=al-qh$ ranges over the nonzero integers $\equiv al\pmod q$; both coordinates in the
minimum are integers, and the infimum is over every fraction $h/l\ne u$, not necessarily reduced. Swapping the coordinates
shows $B(a/q)=B(a^{-1}/q)$ for $q>1$, with the inverse taken modulo $q$; for $q=1$ use the integer $0$ as inverse. In all cases
$B(u)\le1$.

\begin{lemma}[Determinant estimate]\label{lem:determinant}\evidence{Paper proof}
Let $u=a/q$ be reduced, let $z=c/d_z\ne0$ be in lowest terms with $d_z>0$, and put $v=u+z/q^2$. Then
\[
B(v)\ge\min\Bigl\{|z|,\ \frac{B(u)}{2c^2}\Bigr\}.
\]
\end{lemma}

\begin{proof}
Let $h/l\ne v$, put $L=al-qh$ and $w=l/q>0$; then $l^2|v-h/l|=w|L+zw|$. If $L=0$, then $w$ is a positive integer and the value is
at least $|z|$. If $|L|\ge2|z|w$, the value is at least $w|L|/2=l|L|/(2q)\ge B(u)/2$. Otherwise $0<|L|<2|z|w$, and $w=m+(a^{-1}/q)L$ for an
integer $m$. Since $L+zw\ne0$, we have $w+L/z\ne0$, and $(w+Ld_z/c)/L=a^{-1}/q-h'/l'$ with $l'=|cL|$ for a suitable integer $h'$;
this fraction differs from $a^{-1}/q$, so testing $B(a^{-1}/q)=B(u)$ with denominator $|cL|$ gives
$|w+Ld_z/c|\ge B(u)/(c^2|L|)$. Multiplying by $w|z|$ and using $|L|<2|z|w$ gives $w|L+zw|\ge B(u)/(2c^2)$.
\end{proof}

The bound is attained: for $u=1/3$ and $z=1/10$ one has $B(v)=|z|$ through the case $L=0$. Large partial quotients of a new
prefix are thus controlled even when a new prime prevents a literal integral fold of the continued fraction (compare the
folding lemma for $\sum u^{-2^k}$ in \cite{Shallit79,vdPS92}). Since $1/B(u)$ is comparable to the largest partial quotient
of $u$, iterating the lemma bounds the partial quotients of successive prefixes.

\subsection{Anchored upper sets}

A set $U\subseteq\N_0^r$ is \emph{upward closed} if $\mathbf h\in U$ and $\mathbf h'\ge\mathbf h$ coordinatewise imply $\mathbf h'\in U$. Put
$|\mathbf h|=h_1+\cdots+h_r$ and $\mathbf a_r=(\lceil\log_2r\rceil,\ldots,\lceil\log_2r\rceil)$.

\begin{lemma}[Anchored upper sets]\label{lem:upperset}\evidence{Paper proof}
Let $p\ge3$ be prime and $r\ge1$. There is $L=L(r,p)$ such that $\sum_{\mathbf h\in U}p^{|\mathbf h|}\not\equiv0\pmod{p^L}$ for every upward-closed
$U\subseteq\N_0^r$ containing $\mathbf a_r$; the sum converges in $\Z_p$.
\end{lemma}

\begin{proof}
By Dickson's lemma $U$ is the union of the upper orthants of finitely many minimal vectors $\mathbf m_1,\ldots,\mathbf m_M$, $M\ge1$. By
inclusion--exclusion, $\sum_{\mathbf h\in U}x^{|\mathbf h|}=F(x)/(1-x)^r$ with $F(x)=\sum_{\emptyset\ne S}(-1)^{|S|+1}x^{|\vee_{i\in S}\mathbf m_i|}\in\Z[x]$, where
$\vee$ is the coordinatewise maximum, and $F(1)=\sum_{\emptyset\ne S}(-1)^{|S|+1}=1$. Hence $F(p)\equiv1\pmod{p-1}$, so $F(p)\ne0$ because
$p-1\ge2$; evaluation at $x=p$ is valid in $\Z_p$ because $1-p$ is a unit. So each sum is nonzero in $\Z_p$. The anchored upward-closed
sets form a closed subset of the compact space $\{0,1\}^{\N_0^r}$, and the residue of the sum modulo $p^L$ depends only on the
coordinates with $|\mathbf h|<L$; so the sets $Z_L$ of anchored upper sets with sum $\equiv0\pmod{p^L}$ are closed, decrease in $L$, and
have empty intersection. By compactness some $Z_L$ is empty.
\end{proof}

Testing finite truncated upper sets gives a terminating search for $L(r,p)$.

\subsection{The doubling schedule}

\begin{theorem}[Doubling schedule]\label{thm:doubling}\evidence{Paper proof}
Let $b\ge2$, let $p\nmid b$ be an odd prime, let $\mu_k=\ell_j$ at the square indices $k=j^2$ and $\mu_k=1$ otherwise, and let
$u_k$ be the least positive integer with
\[
u_k\equiv\mu_kb^{2^k}\pmod{p^k},\qquad u_k\equiv1\pmod{\rad(b)\mu_k}.
\]
Put $c_k=u_k/(p^k\mu_k)$ and $W=\sum_{k\ge1}c_kb^{-2^k}$. Put $w_k=2^{k-\lceil\sqrt k\rceil}$, and let $W'$ be obtained from $W$ by adding a
coefficient $1$ at every position $2^k+o$ with $o$ odd in $[1,w_k]$ and at $2^k+w_k$, for every $k\ge1$. Then:
\begin{enumerate}[label=\textup{(\roman*)},nosep]
\item $W$ and $W'$ are all-prime: for $2^k\le n<2^{k+1}$ every shifted prefix has reduced denominator exactly
$D_k=p^k\prod_{j^2\le k}\ell_j$;
\item $\mu(W)=\mu(W')=2$;
\item every nonconstant $P\in\Q[x]$ gives a normal number $P(W)$, with $D_N(P(W))\ll_PN^{-\eta_P}$, and a normal number $P(W')$;
in particular $W$ and $W'$ are transcendental;
\item the consecutive position ratios of $W'$ have $\liminf1$ and $\limsup2$, its positions up to $N$ number $N^{1-o(1)}$, and their
density tends to $0$.
\end{enumerate}
\end{theorem}

\begin{proof}
Put $w_k=0$ for $W$ and $m_k=2^k+w_k$ in both cases. \emph{Denominators.} The moduli $p^k$ and $\rad(b)\mu_k$ are coprime, so
$0<u_k\le p^k\rad(b)\mu_k$ and $0<c_k\le\rad(b)$; $u_k$ is a unit at $p$ and at every prime of $b\mu_k$. The last distinguished term
has the unique deepest $p$-pole $p^{-k}$, each auxiliary prime occurs in exactly one coefficient with a unit numerator, and the
integer fillers change no denominator; this gives (i), and every prime $\ell\nmid b$ eventually divides $D_k$. At $m_k$ the
unshifted prefix $U_k$ has exact denominator $q_k=D_kb^{m_k}$: after multiplication by $D_kb^{m_k}$ every earlier term carries a
factor $b$, while the last term is a unit at each prime of $b$. Also $\log D_k=O(k+\sqrt k\log(k+1))$.

\emph{Exponent.} Write $U_{k+1}=U_k+z_k/q_k^2$. The next cluster sums to $A/(p^{k+1}\mu_{k+1}b^{2^{k+1}+w_{k+1}})$ with a positive integer
$A\ll_bp^{k+1}\mu_{k+1}b^{w_{k+1}}$, by the geometric sum of the coefficient bounds. Since $2\cdot2^k=2^{k+1}$,
$z_k=AD_k^2b^{2w_k-w_{k+1}}/(p^{k+1}\mu_{k+1})$, and the numerator and denominator of the reduced $z_k$ have logarithms
$O_b(k+\log D_k+\log\mu_{k+1}+w_k+w_{k+1})=o(2^k)$; this also bounds $\log^+(1/|z_k|)$. By \Cref{lem:determinant},
$\log(1/B(U_{k+1}))\le\max(\log(1/|z_k|),\log(1/B(U_k))+2\log|c_{z_k}|+\log2)$, where $c_{z_k}$ is the numerator of $z_k$; summing a
sequence of terms $o(2^i)$ up to $k$ gives $\log(1/B(U_k))=o(2^k)$. Given a large $l$, choose the first $k$ with $2^k\ge4\log_bl$. Then
$2^k<8\log_bl$, $q_k\ge b^{2^k}\ge l^4>l$, so $h/l\ne U_k$ for every $h$, and
\[
\Bigl|U_k-\frac hl\Bigr|\ge\frac{B(U_k)}{l^2}=l^{-2-o(1)},\qquad 0<W^{(\prime)}-U_k\ll_bb^{-2^{k+1}}\le l^{-8},
\]
uniformly in $h$, where $W^{(\prime)}$ is the chosen series. Hence $|W^{(\prime)}-h/l|\ge l^{-2-o(1)}$; this excludes rationality (test multiples
of a hypothetical fixed denominator) and gives $\mu\le2$, and Dirichlet's theorem gives $\mu\ge2$.

\emph{Pole retention in degree $r$.} For $2^K\le n<2^{K+1}$ reflect distinguished indices by $h_i=K-i$. An ordered $r$-tuple of
distinguished indices is present at $n$ exactly when its reflection lies in the upward-closed set
$U=\{\mathbf h\in\N_0^r:\sum_i2^{-h_i}\le n/2^K\}$, which contains $\mathbf a_r$ because $r2^{-\lceil\log_2r\rceil}\le1$. Let $L=L(r,p)$ be as in
\Cref{lem:upperset}. Normalize the unshifted prefix of the $r$-th power by $p^{rK}$. A tuple of loss $|\mathbf h|\ge L$ vanishes modulo
$p^L$; for $|\mathbf h|<L$ all indices satisfy $i\ge K-L+1$, and for $K\ge2L$ they are valid indices with
$u_ib^{-2^i}/\mu_i\equiv1\pmod{p^L}$ by construction. So the surviving mass is the upper-set sum modulo $p^L$, which is nonzero.
Every tuple using a filler, and every term of lower degree, has depth at most $(r-1)K$. Consequently, for every rational
polynomial of degree $r$ the shifted prefix has $p$-exponent $rK+O_P(1)$ at \emph{every} sufficiently late shift, also when
positions of different tuples coincide, since the entire normalized prefix is evaluated as one upper-set sum.

\emph{Digits.} Clear denominators as in \Cref{thm:allprime}. For $K=\lfloor\log_2N\rfloor$, all insertions through $N$ lie in the
$O_r((K+1)^r)$ intervals $[v,v+rw_K]$ about sums $v$ of at most $r$ distinguished positions (including $v=0$). Remove these
intervals, the shifts $n<\sqrt N$, buffers of length $C_P\log N$ before each interval and before $N$, and the remaining components
shorter than $N^{1/3}$; this costs $O_P(\sqrt N+(K+1)^r(1+rw_K+\log N+N^{1/3}))=o(N)$, a power saving for $W$. The coefficients grow
polynomially in the exponent, so off the buffers the real tail is $O(N^{-1})$. On a retained run, starting at or above $\sqrt N$
(so its shell index is at least $\frac12\log_2N$) and of length $T'\ge N^{1/3}$, the complementary denominator divides a fixed integer
times $\prod_{j^2\le K}\ell_j^r$, whose logarithm is $O(\sqrt K\log(K+1))$; so the stride is $\tau=N^{o(1)}$, $M_0=\lceil T'/\tau\rceil\ge N^{1/3-o(1)}$,
and $A(b^\tau)\le a_0+\log_p\tau=o(\log N)$. For a small fixed $\xi>0$ the pole exponent $e$ after multiplication by any $h\le N^\xi$
satisfies $c_1\log N\le e\le c_2\log N$, so $p^{\theta e}\le M_0$ for a fixed $\theta>0$. \Cref{lem:uniformorbit}, \Cref{cor:split} and the
Erd\H{o}s--Tur\'an inequality with $N^\xi$ frequencies give a power saving on each retained run, and \Cref{lem:blocks} proves
normality, with the power rate for $W$. Normality of every nonconstant $P(W^{(\prime)})$ excludes algebraic relations.

(iv) For $2^k\le N<2^{k+1}$ the preceding complete cluster alone has $\gg2^{k-1-\lceil\sqrt{k-1}\rceil}=N^{1-o(1)}$ positions, while
$\sum_{j\le k}w_j=o(2^k)$; within a cluster the consecutive positions differ by at most $2$, and the ratio $2^{k+1}/(2^k+w_k)\to2$.
\end{proof}

The proof permits arbitrary subsets of integer offsets in $[1,w_k]$ containing the endpoint; the stated choice gives a concrete
count. For $W'$ the removed cluster windows cost $O(K^r2^{K-\sqrt K})=o(N)$, which is not a power saving.

\subsection{A positive series with consecutive ratios tending to one}

\begin{theorem}[Ratio one]\label{thm:dense}\evidence{Paper proof}
Keep $c_k$ and the distinguished terms $c_kb^{-2^k}$ of \Cref{thm:doubling}. Let $h_k$ be the largest power of two at most $\sqrt k$,
and in each shell $(2^k,2^{k+1}]$ add a coefficient $1$ at the positions $2^k+jh_k$, $1\le j\le2^k/h_k$, combining the endpoint $2^{k+1}$ with
the next distinguished term. Call the resulting positive series $V$, and let $n_1<n_2<\cdots$ be its occupied positions. Then
\[
n_j\asymp j\sqrt{\log j},\qquad n_{j+1}/n_j\to1;
\]
$V$ is all-prime; $\mu(V)=2$; and $D_N(P(V))\ll_PN^{-\eta_P}$ for every nonconstant $P\in\Q[x]$.
\end{theorem}

\begin{proof}
\emph{Positions.} Through $N\in[2^k,2^{k+1})$ the preceding complete shell alone gives $\gg2^k/\sqrt k$ positions, and summing $2^i/h_i$ up to
$k$, split at $k/2$, gives the matching upper bound; so $j\asymp n/\sqrt{\log n}$ at $n=n_j$. The largest gap in the shell is
$h_k\le\sqrt k$, so $n_{j+1}/n_j\to1$. Integer insertions preserve the exact shifted denominators $D_k$, so $V$ is all-prime.

\emph{Regrouping.} Put $\lambda_0=0$ and $\lambda_k=1/(b^{h_k}-1)$. Since $h_k\mid2^k$, finite geometric summation gives
\[
\sum_{j=1}^{2^k/h_k}b^{-2^k-jh_k}=\lambda_k\bigl(b^{-2^k}-b^{-2^{k+1}}\bigr),\qquad V=\sum_{k\ge1}d_kb^{-2^k},\quad d_k=c_k+\lambda_k-\lambda_{k-1},
\]
an absolutely convergent signed expansion with bounded coefficients. The orbit $\fr{b^nP(V)}$ is a property of the real number,
and we compute it from this second expansion, whose positions are the doubling positions; the dense raw support is never used.

\emph{$p$-adic control.} Since $h_{k-1}\mid h_k$, both correction denominators divide $b^{h_k}-1$. If $t\mid h_k$, lifting the exponent
gives $v_p(b^{h_k}-1)=a_0+v_p(h_k/t)=a_0$, because $h_k/t$ is a power of $2$; otherwise $v_p(b^{h_k}-1)=0$. So in a reflected tuple the
normalized unit $p^ib^{-2^i}d_i$ is $\equiv1\pmod{p^{i-a_0}}$ for $i>a_0$, and early indices have fixed finite depth. For a late shell and
fixed precision $L$, tuples using an early index have depth defect at least $K-O(1)$ and vanish after normalization by $p^{rK}$;
all others vanish for loss at least $L$ and agree with the upper-set mass for loss below $L$ once $K\ge2L+a_0$. The complementary
denominator of a degree-$r$ prefix divides, up to the fixed polynomial denominator, the prime-to-$bp$ part of
$\bigl((b^{h_K}-1)\prod_{j^2\le K}\ell_j\bigr)^r$, whose logarithm is $O_{P,b}(\sqrt K\log(K+1))=o(K)$. The regrouped coefficients are bounded,
and there are $O_P((1+\log N)^r)$ insertions through $N$; the proof of \Cref{thm:doubling} for $W$ applies, including the power
saving. The period $h_K$ of the filler is absorbed in the stride of \Cref{cor:split}.

\emph{Exponent.} Let $\widetilde U_k=\sum_{i\le k}d_ib^{-2^i}$ and $q_k=\den\widetilde U_k$. A common off-radix denominator is
$E_k=D_k(b^{h_k}-1)$, with $\log E_k=O_b(k+\sqrt k\log(k+1))$. Put $A_k=E_kd_k$, an integer with $0<|A_k|\ll_bE_k$ for $k>a_0$, since
$v_p(d_k)=-k$ then. The numerator of $\widetilde U_k$ over $E_kb^{2^k}$ is $\equiv A_k\pmod{b^{2^{k-1}}}$, and $|A_k|<b^{2^{k-1}}$ for large $k$, so
each prime of $b$ has the same valuation in that numerator as in $A_k$. Consequently $q_k/b^{2^k}$ is a rational whose numerator
and denominator are at most $\exp(O_b(k+\sqrt k\log(k+1)))$. In
\[
\widetilde U_{k+1}-\widetilde U_k=\frac{z_k}{q_k^2},\qquad z_k=d_{k+1}\Bigl(\frac{q_k}{b^{2^k}}\Bigr)^2,
\]
both heights of the nonzero signed correction have logarithm $O_b(k+\sqrt k\log(k+1))$. Applying \Cref{lem:determinant} from a
fixed index gives $\log(1/B(\widetilde U_k))=O_b(k^2+k^{3/2}\log(k+1))=o(2^k)$, and the absolute tail is $O_b(b^{-2^{k+1}})$. Choosing
$2^k\in[4\log_bl,8\log_bl)$ gives $q_k\ge b^{2^k-O(k+\sqrt k\log(k+1))}>l$ and $|V-h/l|\ge l^{-2-o(1)}$, which proves $\mu(V)=2$.
\end{proof}

A raw positive prefix ending at an ordinary insertion $n_j$ has exact radix factor $b^{n_j}$, since its last coefficient is $1$,
and at $n_j=2^k$ the same numerator argument gives $\log_b\den=n_j+o(n_j)$; the positive raw tail begins at $n_{j+1}$ with
coefficient at least $1$, so its negative base-$b$ logarithm is $n_{j+1}+O_b(1)$. Thus \eqref{eq:gapcondition} genuinely fails on the
raw prefixes of $V$, and \Cref{prop:certificate} certifies nothing for them, while the regrouped prefixes prove the optimal
exponent. The distinction is between the obstruction to a certificate from a given sequence of prefixes and the irrationality
exponent of the limit.

\section{Questions}\label{sec:questions}

\begin{question}\label{q:ratioone}
Is the irrationality exponent of the series $W_a$ of \Cref{thm:hybrid}, or of the series of \Cref{rem:quasipoly}, finite?
Is either a Liouville number? (Their coefficients admit no geometric regrouping of the kind used in \Cref{thm:dense}.)
\end{question}

\begin{question}
Which irrationality exponents in $(2,\infty)$ occur for all-prime series all of whose nonconstant rational polynomial images are
normal? \Cref{thm:doubling} gives $2$ and \Cref{thm:stoneham} gives every odd prime; \Cref{rem:realratio} gives every real number
at least $(3+\sqrt5)/2$, but only in degree at most two.
\end{question}

\begin{question}
How early can the auxiliary primes appear? In every construction here of degree at least two the $j$-th prime first appears at a
position $n_j$ with $\log n_j\gg j\log j$. Does a version of \Cref{thm:geometric} with a new prime at every index exist? A
short-orbit estimate uniform over growing auxiliary support would be the natural input.
\end{question}

\begin{question}
Does $P(W')$ in \Cref{thm:doubling} admit a power saving? Is there an all-prime series with $\mu=2$ all of whose polynomial images
are normal and whose positions have positive density?
\end{question}

\section{Formal verification and data availability}\label{sec:formal}

The Lean~4 \cite{Lean4} declarations cited above, built with Mathlib \cite{Mathlib}, are in the repository
\url{https://github.com/eldernos/MathematicsResearch}. Links to \ttt{SurvivingGrid} and \ttt{PrimePowerGrid} are pinned to the
commit
\begin{center}\ttt{\phaseXIXcommit},\end{center}
and links to \ttt{ProtectedPole} to the commit
\begin{center}\ttt{\phaseXXIcommit}.\end{center}
At each commit the full build completed and the axiom audit accepted only \ttt{propext}, \ttt{Classical.choice} and
\ttt{Quot.sound}, and an independent kernel replay covered the modules used here. The Lean declarations used in
\Cref{lem:grid,lem:poledom,thm:allprime,thm:protected} are finite arithmetic cores: the order of a principal unit and the count
over one stretch of a grid, the exact $p$-adic depth of the prefixes of $W$, the divisibility of their denominators by the newest
auxiliary prime, and the divisibility lower bound for a truncated Taylor expression. None of them states an exponential-sum
estimate, a normality conclusion, a rate or a Diophantine property; no statement of this paper is kernel-verified as a whole.
All other statements have written proofs; \Cref{thm:freeze,thm:wsmallball} are quoted from \cite{PaperVIIv2}.

Exact rational computations exercising these statements are in the repository:
\href{https://github.com/eldernos/MathematicsResearch/tree/\phaseXXIcommit/scripts/phase20_normality.py}{\ttt{scripts/phase20\_normality.py}}
(exact prefix denominators of $W$ in \Cref{thm:allprime} at 270 prefixes and five radix/prime pairs) and
\href{https://github.com/eldernos/MathematicsResearch/tree/\phaseXXIcommit/scripts/phase21_nonlinear.py}{\ttt{scripts/phase21\_nonlinear.py}}
at the commit above; \ttt{scripts/phase22\_short\_orbits.py} (\Cref{lem:uniformorbit,cor:split,lem:exactorder}, the exact depth of
\Cref{thm:geometric} at $145$ insertion positions, and the Stoneham congruences of \Cref{thm:stoneham}); and
\ttt{scripts/phase23\_nonlinear\_short\_orbits.py} (\Cref{lem:determinant} on $1960$ signed corrections, \Cref{lem:upperset} on $251$
antichains, the exact denominators, correction heights and partial quotients of \Cref{thm:doubling,thm:dense}, the regrouping of
\Cref{thm:dense}, direct convolutions of its squares and cubes, and the exact auxiliary orders of \Cref{lem:exactorder}), which
accompany this draft in the same repository. They are finite checks and illustrations, not proofs of the asymptotic statements. The
theorem verification map and the novelty ledger are the files \ttt{VERIFICATION.md} and \ttt{LEDGER.md} next to the source of this
paper (directory \ttt{papers/12-all-prime-normality-exponents}).

\section{Use of AI tools}

This manuscript was produced in an AI-assisted research project. Large language models proposed conjectures and proof
strategies, wrote the proofs and the text, implemented the exact computations and Lean~4 formalizations, searched the
literature and reviewed the source notes across several model providers. The repository owner directed the project and is
responsible for any release. The Lean kernel checked the declarations listed in the verification map. The written proofs,
literature comparisons and priority statements have been checked only by AI models from different providers: each source note
received one cross-provider model review, whose findings were repaired and re-read, and this assembled manuscript has not been
reviewed as a whole. No independent human expert has checked the proofs. Cross-model review is not independent human peer review.

\bibliographystyle{plainurl}
\bibliography{references}
\end{document}
